\section{\texorpdfstring{The Tower over $\Q(\sqrt{241})$}{The Tower over Q(sqrt 241)}}\label{tw:tower}

This section constructs the number fields of Theorem~\ref{tw:field-family}
as finite subextensions of an infinite pro-$2$ extension of $B=\Q(\sqrt{241})$
with prescribed local Galois groups: ramification is allowed only above
$2$, $3$ and $5$, and the order of the Frobenius is bounded by four at the
two primes $\mathfrak t_1,\mathfrak t_2$ above $29$ and at one of the two
primes $\mathfrak t_3,\mathfrak t_4$ above $41$, namely $\mathfrak t_3$
(Lemma~\ref{tw:base}). Subsection~\ref{tw:base-section} describes $B$ and its
$\{2,3,5\}$-units. Subsection~\ref{tw:local-section} recalls the local Galois
groups involved and constructs the dyadic local field.
Subsection~\ref{tw:global-section} shows that the Galois group
$G_S$ of the maximal pro-$2$ extension of $B$ unramified outside
$2,3,5$ has a presentation with eight generators and seven local relators
(Lemma~\ref{tw:complete-global-presentation}).
Subsection~\ref{tw:prescribed-section} imposes the local conditions and
defines the quotient $G_B$ of $G_S$ (Definition~\ref{tw:GB}).
Subsection~\ref{tw:retained-section} computes the first three graded pieces
of the Zassenhaus filtration $(D_nG_B)_n$ of $G_B$, recalled below: the
prescribed local groups inject into $G_B/D_3G_B$, and the complex
conjugation $\iota_1$ (Definition~\ref{tw:local-generators}) has $2^{15}$
conjugates in $G_B/D_4G_B$ (Lemma~\ref{tw:retained-quotient}).
Subsections~\ref{tw:fox-section} and~\ref{tw:infinite-section} prove that
$G_B$ is infinite, by a filtered Fox calculus and the Golod--Shafarevich
method \cite{GolodShafarevich}: the Golod--Shafarevich function $P_B$ of
\eqref{tw:PB} takes a negative value (Lemma~\ref{tw:PB-value}).
Subsection~\ref{tw:fields-section} proves Theorem~\ref{tw:field-family}.
Later sections use Theorem~\ref{tw:field-family} (the fields, their root
discriminant, signature and local types), the elementary images of
Table~\ref{tw:vector-table} (Section~\ref{gf:section} and the census of
Section~\ref{an:section}), and the description
$\operatorname{gr}_2G_B=\mathfrak L_2/R_2$ of Lemmas~\ref{tw:graded}
and~\ref{tw:quadratic-layer} (Sections~\ref{an:section}
and~\ref{dh:section}).

We use the following conventions. Subgroups of profinite groups are closed,
generation is topological, and the normal closure of a set is the smallest
closed normal subgroup containing it. The commutator is
$[g,h]=g^{-1}h^{-1}gh$. The \emph{generator rank} and the \emph{relation
rank} of a pro-$2$ group $G$ are $\dim H^1(G,\Ftwo)$ and $\dim H^2(G,\Ftwo)$.
For an extension $M/A$ of local or global fields, $\mathfrak D_{M/A}$ denotes
its different and $\mathfrak d_{M/A}$ its relative discriminant. For a finite
$2$-group $\Delta$ with augmentation
ideal $\mathfrak I\subset\Ftwo[\Delta]$, the Zassenhaus filtration is
$D_n\Delta=\{g:g-1\in\mathfrak I^n\}$; for a pro-$2$ group it is defined
through finite quotients. Then $[D_m,D_n]\subset D_{m+n}$, $g^2\in D_{2n}$
for $g\in D_n$, a surjection maps $D_n$ onto $D_n$, and
$\operatorname{gr}\Delta=\bigoplus_{n\ge1}\operatorname{gr}_n\Delta$, with
graded pieces $\operatorname{gr}_n\Delta=D_n\Delta/D_{n+1}\Delta$, is a
restricted Lie algebra over $\Ftwo$ whose bracket is induced by commutators
and whose restricted square $v\mapsto v^{[2]}$ is induced by squaring
\cite{Jennings,Quillen,Ershov2012}. By Lazard's formula
\cite[Proposition~3.2]{Ershov2012}, $D_n=\prod_{i2^j\ge n}\gamma_i^{2^j}$,
where $\gamma_i$ is the lower central series. In particular $D_1$ is the
whole group, $D_2$ is the Frattini subgroup, and $D_3=D_2^{\,2}[D_2,D_1]$. For
a free pro-$2$ group $\mathfrak F$ of rank $n$ (free pro-$2$ groups are
written in fraktur: $\mathfrak F$, $\mathfrak F_k$, $\mathfrak F'$),
$\operatorname{gr}\mathfrak F$ is the free restricted Lie algebra on $\mathfrak F/D_2\mathfrak F\cong\Ftwo^n$
\cite{Quillen,Ershov2012}. We realize it inside the free associative
algebra $\Ftwo\langle X_1,\dots,X_n\rangle$, with bracket $[a,b]=ab+ba$ and
restricted square $a^{[2]}=a^2$; its graded pieces of degrees one, two and
three have dimensions $n$, $n(n+1)/2$ and $(n^3-n)/3$.

\begin{definition}\label{tw:hilbert-def}
A \emph{filtered space} is a finite-dimensional $\Ftwo$-vector space $M$
with subspaces
$M=\operatorname{Fil}^0M\supseteq\operatorname{Fil}^1M\supseteq\cdots$ such
that $\operatorname{Fil}^mM=0$ for large $m$. Its \emph{Hilbert polynomial}
is
$h_M(t)=\sum_{m\ge0}\dim(\operatorname{Fil}^mM/\operatorname{Fil}^{m+1}M)\,t^m$.
For a finite $2$-group $\Delta$ we filter $\Ftwo[\Delta]$ by the powers
$\mathfrak I^m$ of its augmentation ideal $\mathfrak I$, so that
$h_{\Ftwo[\Delta]}(t)=\sum_{m\ge0}\dim(\mathfrak I^m/\mathfrak I^{m+1})\,t^m$.
\end{definition}

\begin{theorem}[Jennings \cite{Jennings}; see also \cite{Quillen}]\label{tw:jennings}
Let $\Delta$ be a finite $2$-group, and let $g_1,\dots,g_m\in\Delta$ be
elements whose classes form a basis of $\operatorname{gr}\Delta$ consisting of
homogeneous elements, the class of $g_i$ lying in
$\operatorname{gr}_{n_i}\Delta$. Then the products
$(g_1-1)^{a_1}\cdots(g_m-1)^{a_m}$ with $a_i\in\{0,1\}$, taken in this order,
form a basis of $\Ftwo[\Delta]$, and those with $\sum_ia_in_i\ge n$ form a
basis of $\mathfrak I^n$. The same holds with the factors taken in any other
fixed order (Remark~\fref{tw:jennings-order}). Consequently
\[
 h_{\Ftwo[\Delta]}(t)=\prod_{n\ge1}(1+t^n)^{\dim\operatorname{gr}_n\Delta}.
\]
\end{theorem}

\begin{remark}\label{tw:jennings-order}
The form of Theorem~\ref{tw:jennings} with the factors in an arbitrary fixed
order follows from Quillen's description of $\operatorname{gr}\Ftwo[\Delta]$
as the restricted enveloping algebra of $\operatorname{gr}\Delta$
\cite{Quillen} and the restricted Poincar\'e--Birkhoff--Witt theorem. We use
this freedom in Lemma~\ref{tw:filtered-local-freeness}(a), where the elements
of a subgroup are placed last.
\end{remark}

\subsection{The Base Field}\label{tw:base-section}
This subsection fixes the $\{2,3,5\}$-units $\alpha_0,\dots,\alpha_7$ of
$B$, the primes above $2,3,5,29,41$ and their square classes
(Lemma~\ref{tw:base}, Table~\ref{tw:squareclass-table}); every later
Kummer-theoretic computation reads off these data.

Let $B=\Q(\sqrt{241})$, with real places $v_1$ ($\sqrt{241}\mapsto+\sqrt{241}$)
and $v_2$ ($\sqrt{241}\mapsto-\sqrt{241}$). Put
\begin{equation}\label{tw:alpha}
\begin{gathered}
 \alpha_0=-1,\qquad \alpha_1=\varepsilon=-71011068+4574225\sqrt{241},\\
 \alpha_2=\tfrac{-6101-393\sqrt{241}}2,\qquad
 \alpha_3=\tfrac{6101-393\sqrt{241}}2,\\
 \alpha_4=31-2\sqrt{241},\quad \alpha_5=31+2\sqrt{241},\quad
 \alpha_6=326-21\sqrt{241},\quad \alpha_7=326+21\sqrt{241}.
\end{gathered}
\end{equation}
Their norms $\mathrm N_{B/\Q}(\alpha_i)$ are $1,-1,-2,-2,-3,-3,-5,-5$, and
$\alpha_2\alpha_3=2$, $\alpha_4\alpha_5=-3$, $\alpha_6\alpha_7=-5$. For a
nonzero ideal $\mathfrak a$ of $\mathcal O_B$ we write
$\mathrm N\mathfrak a=|\mathcal O_B/\mathfrak a|$ for its absolute norm.

\begin{lemma}\label{tw:base}
\begin{enumerate}
\item[(a)] $\mathcal O_B=\Z[(1+\sqrt{241})/2]$, the discriminant of $B$ is
 $241$, and $B$ has class number one.
\item[(b)] The unit $\varepsilon$ has norm $\mathrm N_{B/\Q}(\varepsilon)=-1$, and the
 classes of $-1$ and $\varepsilon$ form a basis of
 $\mathcal O_B^\times/\mathcal O_B^{\times2}$.
\item[(c)] In $\mathcal O_B$ we have $2=\mathfrak p_1\mathfrak p_2$,
 $3=\mathfrak q_1\mathfrak q_2$, $5=\mathfrak r_1\mathfrak r_2$ and
 $29=\mathfrak t_1\mathfrak t_2$, where
 \[
  \mathfrak p_j=\alpha_{1+j}\mathcal O_B,\quad
  \mathfrak q_j=\alpha_{3+j}\mathcal O_B,\quad
  \mathfrak r_j=\alpha_{5+j}\mathcal O_B,\quad
  \mathfrak t_j=(-14127\pm910\sqrt{241})\mathcal O_B,
 \]
 with the sign $+$ for $j=1$. These eight primes are distinct and of degree
 one. Likewise $41=\mathfrak t_3\mathfrak t_4$ with
 $\mathfrak t_3=(621+40\sqrt{241})\mathcal O_B$ and
 $\mathfrak t_4=(621-40\sqrt{241})\mathcal O_B$, two distinct primes of degree
 one. The prime $7$ is inert, and we write $\mathfrak t_0=7\mathcal O_B$, so
 that $\mathrm N\mathfrak t_0=49$. The prime $241$ ramifies. At each of the ten
 primes of degree one above $p\in\{2,3,5,29,41\}$ listed here, the completion
 of $B$ is $\Q_p$, and $\sqrt{241}$ maps to the square root of $241$ in $\Z_p$
 given in Table~\ref{tw:squareclass-table}.
\item[(d)] Let $S=\{\mathfrak p_1,\mathfrak p_2,\mathfrak q_1,\mathfrak q_2,\mathfrak r_1,\mathfrak r_2\}$
 and $V=\mathcal O_B[1/30]^\times/\mathcal O_B[1/30]^{\times2}$. The classes
 of $\alpha_0,\dots,\alpha_7$ form a basis of $V$, and $V$ is the subgroup of
 $B^\times/B^{\times2}$ of classes with even valuation at every prime outside
 $S$. Moreover $\Pic(\mathcal O_B[1/30])=0$.
\item[(e)] The element $\alpha_i$ is a square modulo $\mathfrak t_3$ exactly
 for $i=0,1,5,6,7$, and modulo $\mathfrak t_4$ exactly for $i=0,1,4,6,7$.
 The residue field of $\mathfrak t_0$ is $\mathbb F_{49}$, and $\alpha_i$ is a
 square modulo $\mathfrak t_0$ exactly for $i=0,4,5,6,7$.
\end{enumerate}
\end{lemma}

\begin{proof}
(a) Since $241\equiv1\pmod 4$, the ring of integers and the discriminant are
as stated. The Minkowski bound is $\sqrt{241}/2<8$, so every ideal class
contains an integral ideal of norm at most $7$, which is a product of primes
of norm at most $7$. By (c), these are the six primes above $2,3,5$, and they
are principal.

(b) A direct computation gives $71011068^2-241\cdot4574225^2=-1$. By
Dirichlet's theorem $\mathcal O_B^\times=\{\pm1\}\times\eta^{\Z}$ for a
fundamental unit $\eta$, and $\varepsilon=\pm\eta^k$. Since
$\mathrm N_{B/\Q}(\varepsilon)=-1$, the exponent $k$ is odd, so $-1$ and
$\varepsilon$ generate $\mathcal O_B^\times$ modulo squares. Neither $-1$ nor
$\pm\varepsilon$ is a square, because $B$ is real and
$\mathrm N_{B/\Q}(\pm\varepsilon)=-1$.

(c) We have $241\equiv1\pmod8$, $241\equiv1\pmod3$, $241\equiv1\pmod5$,
$241\equiv9\pmod{29}$, $241\equiv36\pmod{41}$ and $241\equiv3\pmod7$, and $3$
is not a square modulo $7$. Hence $2,3,5,29,41$ split, $7$ is inert and $241$
ramifies. The listed generators have norms $\pm2,\pm3,\pm5,29,41$ (for
instance $621^2-241\cdot40^2=41$), so they generate primes of degree one, and
the displayed products show that the two generators above each $p$ generate
the two distinct primes above $p$; for $29$ and $41$ the product of the two
generators is $p$. A prime of degree one above $p$ is the preimage of
$p\Z_p$ under an embedding $B\to\Q_p$, which sends $\sqrt{241}$ to one of
the two square roots of $241$ in $\Z_p$; the completion there is $\Q_p$, and
the prime contains the stated generator exactly for the square root in
Table~\ref{tw:squareclass-table}.

(d) Every $S$-unit $u$ satisfies
$u\mathcal O_B=\prod_{i\ge2}(\alpha_i\mathcal O_B)^{a_i}$, so
$u\prod_{i\ge2}\alpha_i^{-a_i}$ is a unit, and by (b) the classes of
$\alpha_0,\dots,\alpha_7$ generate $V$. If $\prod_i\alpha_i^{a_i}$ with
$a_i\in\{0,1\}$ is a square, the valuations at $S$ give $a_i=0$ for $i\ge2$,
and then (b) gives $a_0=a_1=0$. If $b\in B^\times$ has even valuation outside
$S$, then
$b\mathcal O_B=\mathfrak a^2\prod_{\mathfrak p\in S}\mathfrak p^{a_{\mathfrak p}}$
with $\mathfrak a$ principal by (a), so $b$ is an $S$-unit times a square.
Finally $\Pic(\mathcal O_B[1/30])$ is a quotient of $\Pic(\mathcal O_B)=0$.

(e) At $\mathfrak t_3$ and $\mathfrak t_4$ the residue field is $\mathbb F_{41}$
and $\sqrt{241}\mapsto6$, respectively $35$, since $621+40\cdot6=21\cdot41$
and $621-40\cdot35=-19\cdot41$; the statement is a computation of Legendre
symbols modulo $41$, repeated by the supplementary programs
\texttt{kummer41.gp} and \texttt{check41.py}. By (c), $\mathrm N\mathfrak t_0=49$. An
element $a\in\mathcal O_B$ prime to
$\mathfrak t_0$ is a square modulo $\mathfrak t_0$ exactly when
$a^{24}\equiv1$, that is, when the reduction of $\mathrm N_{B/\Q}(a)\equiv a^8$ is a
square in $\mathbb F_7$. By the norms listed after \eqref{tw:alpha}, this
holds for $\alpha_i$ exactly when $i=0,4,5,6,7$.
\end{proof}

\begin{table}[htbp]
\centering
\caption{The image of $\sqrt{241}$ in $\Z_p$ at each prime of degree one,
given by its residue, and in $\R$ at the real places; the classes of
$\alpha_0,\dots,\alpha_7$ in $\Q_p^\times/\Q_p^{\times2}$ at these primes, and
their signs at the real places. At $p=2$ the classes are written as
$\pm1,\pm5,\pm2,\pm10$; at odd $p$ as $1,u,p,up$, where $u=-1,2,2,3$ is a
nonsquare unit for $p=3,5,29,41$.}
\label{tw:squareclass-table}
\begin{tabular}{llrrrrrrrr}
\toprule
place & $\sqrt{241}\mapsto$ & $\alpha_0$ & $\alpha_1$ & $\alpha_2$ & $\alpha_3$ & $\alpha_4$ & $\alpha_5$ & $\alpha_6$ & $\alpha_7$\\
\midrule
$\mathfrak p_1$ & $7\bmod 32$ & $-1$ & $-5$ & $-10$ & $-5$ & $1$ & $5$ & $-5$ & $1$\\
$\mathfrak p_2$ & $25\bmod 32$ & $-1$ & $5$ & $5$ & $10$ & $5$ & $1$ & $1$ & $-5$\\
$\mathfrak q_1$ & $2\bmod 3$ & $-1$ & $1$ & $-1$ & $1$ & $3$ & $-1$ & $-1$ & $-1$\\
$\mathfrak q_2$ & $1\bmod 3$ & $-1$ & $-1$ & $-1$ & $1$ & $-1$ & $3$ & $-1$ & $-1$\\
$\mathfrak r_1$ & $1\bmod 5$ & $1$ & $2$ & $2$ & $1$ & $1$ & $2$ & $10$ & $2$\\
$\mathfrak r_2$ & $4\bmod 5$ & $1$ & $2$ & $1$ & $2$ & $2$ & $1$ & $2$ & $10$\\
$\mathfrak t_1$ & $3\bmod 29$ & $1$ & $1$ & $2$ & $1$ & $1$ & $2$ & $2$ & $2$\\
$\mathfrak t_2$ & $26\bmod 29$ & $1$ & $1$ & $1$ & $2$ & $2$ & $1$ & $2$ & $2$\\
$\mathfrak t_3$ & $6\bmod 41$ & $1$ & $1$ & $3$ & $3$ & $3$ & $1$ & $1$ & $1$\\
$\mathfrak t_4$ & $35\bmod 41$ & $1$ & $1$ & $3$ & $3$ & $1$ & $3$ & $1$ & $1$\\
$v_1$ & $+\sqrt{241}$ & $-$ & $+$ & $-$ & $-$ & $-$ & $+$ & $-$ & $+$\\
$v_2$ & $-\sqrt{241}$ & $-$ & $-$ & $+$ & $+$ & $+$ & $-$ & $+$ & $-$\\
\bottomrule
\end{tabular}
\end{table}

The classes in Table~\ref{tw:squareclass-table} follow from the residues of
$\sqrt{241}$ given there; at $p=2$ they need its image modulo $32$. For
example, at $\mathfrak p_1$ we have $4574225\equiv1$ and $-71011068\equiv4$
modulo $8$, so $\varepsilon\equiv4+7\equiv3\equiv-5\pmod 8$. The classes
modulo the inert prime $\mathfrak t_0$ are given by Lemma~\ref{tw:base}(e);
the prime $\mathfrak t_0$ carries no condition in this paper.

\subsection{Local Galois Groups}\label{tw:local-section}
This subsection gives presentations of the local pro-$2$ Galois groups, with
the quadratic initials of their relators (Lemma~\ref{tw:local-groups}), and
constructs the dyadic local field $L_2$ with Galois group $\mathcal D$
(Lemma~\ref{tw:local-two}). These are the local conditions of the tower and
the local inputs of the Golod--Shafarevich estimate of
Subsections~\ref{tw:fox-section} and~\ref{tw:infinite-section}.

For a prime $p$ let $\mathcal G_{\Q_p}$ be the Galois group of the maximal
pro-$2$ extension of $\Q_p$, and for a place $\nu$ of $B$ let $\mathcal G_\nu$
be the Galois group of the maximal pro-$2$ extension of the completion
$B_\nu$; for a real place, $\mathcal G_\nu=\Gal(\C/\R)$. We write
$(\cdot,\cdot)$ for the Hilbert symbol
of order two of a local field \cite[Chapter~III, Section~4]{MilneCFT}. For
odd $p$ and $p$-adic units $u,w$, and for $2$-adic units $u,w$,
\begin{equation}\label{tw:hilbert}
 (p^au,p^bw)_p=(-1)^{ab(p-1)/2}\Bigl(\frac up\Bigr)^b\Bigl(\frac wp\Bigr)^a,
 \qquad
 (2^au,2^bw)_2=(-1)^{\frac{u-1}2\frac{w-1}2+a\frac{w^2-1}8+b\frac{u^2-1}8}
\end{equation}
\cite[Chapter~III, Theorem~1]{SerreCourse}, \cite[Chapter~VIII,
Section~5]{MilneCFT}.

\begin{definition}\label{tw:relator}
Let $g_1,\dots,g_k$ generate a pro-$2$ group $G$, and let $\mathfrak F_k$ be
the free pro-$2$ group on $k$ letters, mapped onto $G$ by sending the letters
to $g_1,\dots,g_k$. A \emph{defining relator} of $G$ with respect to
$g_1,\dots,g_k$ is an element of $\mathfrak F_k$ whose normal closure is the
kernel of $\mathfrak F_k\to G$. If a defining relator $r$ lies in
$D_2\mathfrak F_k$, its \emph{quadratic initial} is its class in
$\operatorname{gr}_2\mathfrak F_k$.
\end{definition}

\begin{lemma}\label{tw:local-groups}
\begin{enumerate}
\item[(a)] If $\nu$ is real, then $\mathcal G_\nu=\langle\iota\rangle\cong C_2$
 with the single defining relator $\iota^2$, and $\iota(\sqrt b)=-\sqrt b$
 exactly when $b<0$.
\item[(b)] Let $p$ be an odd prime and $\varpi$ a uniformizer of $\Q_p$. Then
 $\mathcal G_{\Q_p}$ is generated by a generator $\tau$ of its inertia group
 and a Frobenius lift $\varphi$ with $\varphi(\sqrt\varpi)=\sqrt\varpi$, with
 the single defining relator $\varphi\tau\varphi^{-1}\tau^{-p}$. For
 $b\in\Q_p^\times$, $\tau(\sqrt b)=-\sqrt b$ exactly when $\operatorname{ord}_p(b)$ is odd;
 for a unit $b$, $\varphi(\sqrt b)=-\sqrt b$ exactly when $b$ is not a square
 modulo $p$. The relator has quadratic initial
 $[\tau,\varphi]+\frac{p-1}2\tau^{[2]}$.
\item[(c)] Let $E_2=\Q_2(\sqrt{-1},\sqrt2,\sqrt5)$, and let
 $x,y,z\in\mathcal G_{\Q_2}$ multiply the triple $(\sqrt2,\sqrt{-1},\sqrt5)$
 by the signs $(-,+,+)$, by $(+,-,+)$ and by $(+,-,-)$. Then $x,y,z$
 generate $\mathcal G_{\Q_2}$. If $b\equiv(-1)^{a_1}2^{a_2}5^{a_3}$ modulo
 squares, they multiply $\sqrt b$ by $(-1)^{a_2}$, $(-1)^{a_1}$ and
 $(-1)^{a_1+a_3}$, that is, by the Hilbert symbols $(5,b)_2$, $(-1,b)_2$ and
 $(-2,b)_2$. The group $\mathcal G_{\Q_2}$ has a single defining relator $r$
 with respect to $x,y,z$, and every such $r$ has quadratic initial
 $y^{[2]}+[x,y]+[x,z]$.
\end{enumerate}
\end{lemma}

\begin{proof}
(a) is clear.

(b) By Iwasawa's theorem \cite[Theorem~7.5.3]{NeukirchSchmidtWingberg2008},
the Galois group of the maximal tamely ramified extension of $\Q_p$ is
generated by a Frobenius lift $\varphi'$ and a generator $\tau$ of tame
inertia, with the single relation $\varphi'\tau\varphi'^{-1}=\tau^p$. Since
$p$ is odd,
every $2$-extension of $\Q_p$ is tamely ramified, so $\mathcal G_{\Q_p}$ is
the pro-$2$ group with the same presentation. The extension
$\Q_p(\sqrt\varpi)$ is ramified, so $\tau$ moves $\sqrt\varpi$; replacing
$\varphi'$ by $\varphi'\tau$ if necessary gives $\varphi$. This substitution
does not change the relator, since
$(\varphi'\tau)\tau(\varphi'\tau)^{-1}=\varphi'\tau\varphi'^{-1}$.
The actions on square roots hold because $\Q_p(\sqrt b)$ is ramified exactly
when $\operatorname{ord}_p(b)$ is odd, and because $\varphi$ acts on unramified extensions as
the Frobenius. Modulo $D_3$, the element $\varphi\tau\varphi^{-1}\tau^{-1}$
has class $[\varphi,\tau]=[\tau,\varphi]$, and $\tau^{1-p}$ has class
$\frac{p-1}2\tau^{[2]}$.

(c) We first show that $x,y,z$ generate $\mathcal G_{\Q_2}$, and then read
off the quadratic initial of $r$ from the cup-product pairing on
$H^1(\mathcal G_{\Q_2},\Ftwo)$, which is given by the Hilbert symbol.
By Kummer theory the maximal elementary abelian quotient of
$\mathcal G_{\Q_2}$ is $\Gal(E_2/\Q_2)$, since $\Q_2^\times/\Q_2^{\times2}$
has basis $-1,2,5$. The images of $x,y,z$ form a basis of it, so $x,y,z$
generate \cite[Proposition~3.9.1]{NeukirchSchmidtWingberg2008}. The action on
$\sqrt b$ is read off from the actions on $\sqrt{-1},\sqrt2,\sqrt5$, and it
agrees with the stated Hilbert symbols by \eqref{tw:hilbert}. By
\cite[Theorem~7.5.11]{NeukirchSchmidtWingberg2008}, $\mathcal G_{\Q_2}$ is a
Demu\v skin group of rank $3$: it has a single defining relator $r$, and
$H^2(\mathcal G_{\Q_2},\Ftwo)\cong\Ftwo$. For $b\in\Q_2^\times$ let $\kappa_b$
be the character with $g(\sqrt b)=(-1)^{\kappa_b(g)}\sqrt b$. The characters
dual to $x,y,z$ are $\kappa_2,\kappa_{-5},\kappa_5$. Inflation identifies
$H^2(\mathcal G_{\Q_2},\Ftwo)$ with $H^2$ of the absolute Galois group of
$\Q_2$ \cite[Proposition~7.5.8]{NeukirchSchmidtWingberg2008}, and there the
invariant of $\kappa_a\cup\kappa_b$ is given by the Hilbert symbol $(a,b)_2$
\cite[Chapter~III, Section~4]{MilneCFT}. By \eqref{tw:hilbert}, the matrix of
the indicator of $(a,b)_2=-1$ on $2,-5,5$ is
\[
 \begin{pmatrix}0&1&1\\1&1&0\\1&0&0\end{pmatrix}.
\]
The evaluation $\operatorname{tr}_r$ of the transgression at $r$ is an
isomorphism $H^2(\mathcal G_{\Q_2},\Ftwo)\to\Ftwo$
\cite[Proposition~3.9.12(iii)]{NeukirchSchmidtWingberg2008}, and so is the
invariant, so the two coincide. We now apply the relation--cup identity
\cite[Proposition~3.9.13(ii)]{NeukirchSchmidtWingberg2008}, see also
\cite[Proposition~3.2]{Quadrelli2021}. Write $g_1,g_2,g_3$ for $x,y,z$. The
identity writes the relator as
$r=\prod_jg_j^{qa_j}\prod_{k<l}[g_k,g_l]^{a_{kl}}\,r'$, with $r'$ in the
third term of the descending $q$-central series, and states that
$\operatorname{tr}_r$ of the cup product of the characters dual to $g_k$ and
$g_l$ is $-a_{kl}$ for $k<l$ and $-\binom q2a_k$ for $k=l$. Here $q$, the
smallest elementary divisor of the
abelianization of $\mathcal G_{\Q_2}$, is $2$: by local class field theory
this abelianization is the pro-$2$ completion of $\Q_2^\times$, which is
$\Z_2^2\times\Z/2$. The third term of the $2$-central series is
$(D_2)^2[D_2,D_1]=D_3$ by Lazard's formula, the class of $g_j^2$ in
$\operatorname{gr}_2$ is $g_j^{[2]}$, and the signs disappear modulo $2$,
with $\binom22=1$. Hence the off-diagonal entries of the matrix are the
coefficients of $[x,y],[x,z],[y,z]$ in the quadratic initial of $r$, and the
diagonal entries, which are the values of the cup squares $\kappa\cup\kappa$,
are the coefficients of $x^{[2]},y^{[2]},z^{[2]}$. For $p=2$ the cup square
of a character is its Bockstein, so the diagonal coefficients are also given
by \cite[Proposition~3.9.14]{NeukirchSchmidtWingberg2008}.
\end{proof}

We now construct the dyadic local field. Let $U_4$ be the
unramified extension of $\Q_2$ of degree four, with Frobenius $\mathrm{Fr}$;
let $\mathrm i=\sqrt{-1}$, $k_2=\Q_2(\mathrm i)$, and let
$\pi_5=1+2\mathrm i$, a prime element of $\Z[\mathrm i]$ above $5$. Put
\[
 M_8=\Q_2(\mathrm i,\sqrt5,\sqrt{\pi_5}),\qquad M_{16}=M_8(\sqrt2),\qquad
 L_2=M_{16}U_4 .
\]
Since $\pi_5\bar\pi_5=5$ with $\bar\pi_5=1-2\mathrm i$, the field $M_8$ contains
$\sqrt{\bar\pi_5}=\sqrt5/\sqrt{\pi_5}$. Let $\sigma_y$ and $\sigma_z$ be the
automorphisms of $M_8$ given by
\[
 \sigma_y:\ \mathrm i\mapsto-\mathrm i,\ \sqrt5\mapsto\sqrt5,\ \sqrt{\pi_5}\mapsto\sqrt5/\sqrt{\pi_5};
 \qquad
 \sigma_z:\ \mathrm i\mapsto-\mathrm i,\ \sqrt5\mapsto-\sqrt5,\ \sqrt{\pi_5}\mapsto\sqrt5/\sqrt{\pi_5}.
\]

\begin{lemma}\label{tw:local-two}
\begin{enumerate}
\item[(a)] The extension $M_8/\Q_2$ is Galois of degree $8$, and its Galois
 group $\langle\sigma_y,\sigma_z\rangle$ is dihedral.
\item[(b)] The extension $L_2/\Q_2$ is Galois of degree $32$, with
 ramification index $8$ and residue degree $4$. There are unique
 $x,y,z\in\Gal(L_2/\Q_2)$ such that $x$ is trivial on $M_8$ and $U_4$ and
 negates $\sqrt2$; $y$ restricts to $\sigma_y$ on $M_8$, fixes $\sqrt2$ and
 is trivial on $U_4$; and $z$ restricts to $\sigma_z$ on $M_8$, fixes $\sqrt2$
 and restricts to $\mathrm{Fr}$ on $U_4$. They act on $E_2\subset L_2$ as in
 Lemma~\ref{tw:local-groups}(c), and the map sending the generators of
 \[
  \mathcal D=\langle x,y,z\mid x^2,\ y^2,\ z^4,\ [y,z]^2,\ [x,y],\ [x,z],\ [[y,z],z]\rangle
 \]
 to $x,y,z$ is an isomorphism $\mathcal D\cong\Gal(L_2/\Q_2)$. The inertia group is
 $\langle x,y,[y,z]\rangle\cong C_2^3$.
\item[(c)] The graded pieces of the Zassenhaus filtration of $\mathcal D$ are
 $\operatorname{gr}_1\mathcal D=\langle\bar x,\bar y,\bar z\rangle\cong\Ftwo^3$ and
 $\operatorname{gr}_2\mathcal D=\langle\overline{z^2},\overline{[y,z]}\rangle\cong\Ftwo^2$,
 and $D_3(\mathcal D)=1$, so the Hilbert polynomial of $\Ftwo[\mathcal D]$
 (Definition~\ref{tw:hilbert-def}) is $h_{\Ftwo[\mathcal D]}(t)=(1+t)^3(1+t^2)^2$.
\item[(d)] The different $\mathfrak D_{L_2/\Q_2}$ has valuation $18$ in $L_2$;
 equivalently, $\operatorname{ord}_2(\mathfrak D_{L_2/\Q_2})=18/8=9/4$, where
 $\operatorname{ord}_2$ is the valuation of $L_2$ normalized by
 $\operatorname{ord}_2(2)=1$.
\end{enumerate}
\end{lemma}

\begin{proof}
\emph{Quadratic extensions of $k_2$.} Let $\operatorname{ord}_{k_2}$ be the
normalized valuation of $k_2$. Put $\varpi=1+\mathrm i$, a uniformizer of $k_2$,
so that $\operatorname{ord}_{k_2}(2)=2$, and for $b\in k_2^\times$ let $c(b)$ be the exponent of the
conductor of $k_2(\sqrt b)/k_2$, which for a quadratic extension equals the
exponent of its discriminant. Then $c(5)=0$, since $\Q_2(\sqrt5)/\Q_2$ is
unramified. Since $\pi_5=1+\varpi^2$, the element
$\varpi'=(\sqrt{\pi_5}-1-\varpi)/\varpi$ is a root of the Eisenstein
polynomial $X^2+\frac{2(1+\varpi)}\varpi X+\frac2\varpi$; its derivative at
$\varpi'$ is the sum of terms of valuations $5$ and $2$ in
$k_2(\sqrt{\pi_5})$, so $c(\pi_5)=2$. Next $2\pi_5=\varpi^2u'$ with
$u'=2-\mathrm i$, and $\sqrt{u'}-1$ is a root of $X^2+2X+1-u'$, which
is Eisenstein since $\operatorname{ord}_{k_2}(1-u')=1$; its derivative $2\sqrt{u'}$ has
valuation $4$, so $c(2\pi_5)=4$. Similarly $2=-\mathrm i\varpi^2$, $-\mathrm i\equiv\mathrm i$
modulo squares, and $\sqrt{\mathrm i}-1$ is a root of $X^2+2X+1-\mathrm i$, so $c(2)=4$.
Finally $c(5b)=c(b)$ when $b\notin k_2^{\times2}\cup5k_2^{\times2}$: the field
$k_2(\sqrt5,\sqrt b)$ is unramified over both $k_2(\sqrt b)$ and $k_2(\sqrt{5b})$,
so the two ways of computing its discriminant over $k_2$ give $2c(b)=2c(5b)$.

(a) Since $k_2(\sqrt5)/k_2$ is unramified and $k_2(\sqrt{\pi_5})/k_2$ is ramified,
$\pi_5$ is not a square in $k_2(\sqrt5)$, so $[M_8:\Q_2]=8$. The field $M_8$
contains the conjugate $\sqrt{\bar\pi_5}$ of $\sqrt{\pi_5}$, so it is Galois
over $\Q_2$, and an automorphism is determined by the images $\pm\mathrm i$,
$\pm\sqrt5$ and a square root of the image of $\pi_5$; in particular
$\sigma_y$ and $\sigma_z$ exist. Now $\sigma_y^2=1$, $\sigma_z^2$ fixes
$\mathrm i,\sqrt5$ and negates $\sqrt{\pi_5}$, $\sigma_z$ has order four, and
$\sigma_y\sigma_z\sigma_y=\sigma_z^{-1}$. So
$\Gal(M_8/\Q_2)=\langle\sigma_y,\sigma_z\rangle$ is dihedral of order eight,
with center $\langle\sigma_z^2\rangle$ and maximal abelian subextension
$\Q_2(\mathrm i,\sqrt5)$.

(b) The quadratic subfields of $M_8$ are $\Q_2(\sqrt b)$ for $b=-1,5,-5$, so
$\sqrt2\notin M_8$, and $M_{16}$ has degree $16$ and group
$\Gal(M_8/\Q_2)\times C_2$. This group has elementary abelian
abelianization, so the maximal unramified subextension of $M_{16}$, which is
cyclic over $\Q_2$ and contains $\Q_2(\sqrt5)$, is $\Q_2(\sqrt5)$. Hence
$M_{16}$ has ramification index $8$, $M_{16}\cap U_4=\Q_2(\sqrt5)$, and
$[L_2:\Q_2]=16\cdot4/2=32$. Since $L_2$ contains $U_4$ and has ramification
index at least $8$, its residue degree is $4$ and its ramification index is
$8$. Restriction identifies $\Gal(L_2/\Q_2)$ with the set of triples
$(\sigma,s,\mu)$, where $\sigma\in\Gal(M_8/\Q_2)$, $s=\pm1$ is
the sign by which the element multiplies $\sqrt2$, and $\mu\in\Gal(U_4/\Q_2)$
agrees with $\sigma$ on $\sqrt5$; both sets have $32$ elements. Since $\mathrm{Fr}$
negates $\sqrt5$, the triples
\[
 x=(1,-1,1),\qquad y=(\sigma_y,1,1),\qquad z=(\sigma_z,1,\mathrm{Fr}),\qquad
 w=[y,z]=(\sigma_z^2,1,1)
\]
exist and are the unique elements described in the statement. They act on
$E_2$ as required. The elements $x$ and $w$ are central involutions, $y$ is
an involution, and $z$ has order four. Hence the relations of $\mathcal D$ hold. Conversely, in
$\mathcal D$ these relations make $w=[y,z]$ central: $w^y=w^{-1}=w$, and $[w,z]=1$ is
imposed. Since $zy=yzw$, every element of $\mathcal D$ is
$x^{a_1}y^{a_2}z^{a_3}w^{a_4}$ with $a_1,a_2,a_4\in\Ftwo$ and
$a_3\in\Z/4$, and
\[
 (a_1,a_2,a_3,a_4)(a'_1,a'_2,a'_3,a'_4)
 =(a_1+a'_1,\,a_2+a'_2,\,a_3+a'_3,\,a_4+a'_4+a_3a'_2).
\]
These $32$ normal forms have distinct images in $\Gal(L_2/\Q_2)$: the third
component determines $a_3$, the second $a_1$, and
$\sigma_y^{a_2}\sigma_z^{a_3+2a_4}$ determines $a_2$ and $a_4$. So
$\mathcal D\cong\Gal(L_2/\Q_2)$. The inertia group consists of the elements trivial
on $U_4$, namely $\langle x,y,w\rangle$.

(c) The Frattini subgroup of $\mathcal D$ is $\langle z^2,w\rangle$, which is central
of exponent two, and $\mathcal D$ has exponent four. So Lazard's formula gives
$D_3(\mathcal D)=[[\mathcal D,\mathcal D],\mathcal D]\,[\mathcal D,\mathcal D]^2\mathcal D^4=1$, and the Hilbert polynomial follows from
Theorem~\ref{tw:jennings}.

(d) The extension $M_{16}/k_2$ is abelian with group $C_2^3$, generated by
$\sqrt5,\sqrt{\pi_5},\sqrt2$. We apply the conductor--discriminant formula for
abelian extensions \cite[Chapter~V, Theorem~3.27]{MilneCFT} to the extension
$\Q(\mathrm i)(\sqrt5,\sqrt{1+2\mathrm i},\sqrt2)$ of $\Q(\mathrm i)$, whose completion at the unique
prime above $1+\mathrm i$ is $M_{16}/k_2$. Taking $(1+\mathrm i)$-adic valuations gives
\[
 \operatorname{ord}_{k_2}(\mathfrak d_{M_{16}/k_2})=\sum_bc(b)=0+0+2+2+4+4+4+4=20,
\]
the sum over the classes $b=1,5,\pi_5,5\pi_5,2,10,2\pi_5,10\pi_5$. Hence
\[
 \operatorname{ord}_2(\mathfrak d_{M_{16}/\Q_2})
 =[M_{16}:k_2]\,\operatorname{ord}_2(\mathfrak d_{k_2/\Q_2})+20=8\cdot2+20=36.
\]
Since $M_{16}$ has residue degree two over $\Q_2$, its different has
valuation $36/2=18$. The extension $L_2/M_{16}$ is unramified, so the
different of $L_2/\Q_2$ also has valuation $18$, and $18/8=9/4$.
\end{proof}

\begin{remark}\label{tw:triple-remark}
The field $E_2$ is the maximal elementary abelian subextension of $L_2$, so
$\Gal(L_2/E_2)$ is the Frattini subgroup $\langle z^2,[y,z]\rangle$, which is
central of exponent two. Hence any triple in $\Gal(L_2/\Q_2)$ that acts on
$E_2$ as $x,y,z$ do differs from $x,y,z$ by central elements of order at
most two; it satisfies the relations of $\mathcal D$ and generates, so it also
defines an isomorphism $\mathcal D\cong\Gal(L_2/\Q_2)$. The description of the
inertia group, however, refers to the triple of
Lemma~\ref{tw:local-two}(b).
\end{remark}

\subsection{The Global Group and Its Complete Presentation}\label{tw:global-section}
This subsection introduces the Galois group $G_S$ of the maximal pro-$2$
extension of $B$ unramified outside $2,3,5$, its local generators and their
elementary images (Lemma~\ref{tw:vectors}), and shows that $G_S$ has $8$
generators and $7$ relations, given by local relators
(Lemma~\ref{tw:complete-global-presentation}). The presentation is the
starting point of Subsections~\ref{tw:prescribed-section}--\ref{tw:infinite-section}.

Let $B_S$ be the maximal pro-$2$ extension of $B$, in a fixed algebraic
closure, that is unramified at every finite prime outside $S$; no condition
is imposed at $v_1,v_2$. Let $G_S=\Gal(B_S/B)$, and let
$E=B(\sqrt{\alpha_0},\dots,\sqrt{\alpha_7})=B(\sqrt V)$ be the
\emph{Kummer field}.

\begin{lemma}\label{tw:kummer-field}
The field $E$ is the maximal elementary abelian extension of $B$ contained in
$B_S$, and $[E:B]=256$. Hence $G_S/D_2G_S=\Gal(E/B)$.
\end{lemma}

\begin{proof}
Every prime above $2$ lies in $S$, so a quadratic extension $B(\sqrt b)$ lies
in $B_S$ exactly when $b$ has even valuation outside $S$. By
Lemma~\ref{tw:base}(d), $E$ is the maximal elementary abelian extension of
$B$ in $B_S$, and $[E:B]=256$. Since $D_2G_S$ is the Frattini subgroup of
$G_S$, its fixed field is this maximal elementary abelian extension.
\end{proof}

\begin{definition}\label{tw:elementary-image}
Let $M\subseteq B_S$ be a Galois extension of $B$ containing $E$, for
example $M=B_S$. The \emph{elementary image} of $g\in\Gal(M/B)$ is the vector
$\bar g\in\Ftwo^8$ with $g(\sqrt{\alpha_i})=(-1)^{\bar g_i}\sqrt{\alpha_i}$ for
$0\le i\le7$. For $a\in V$, the \emph{Kummer character} $\kappa_a$ is defined by
$g(\sqrt a)=(-1)^{\kappa_a(g)}\sqrt a$; its values lie in $\Ftwo$. (The
$\pm1$-valued characters $\chi_e$ of Section~\ref{gf:section} are
$\chi_e=(-1)^{\kappa_{a}}$ with $a=\prod_i\alpha_i^{e_i}$.)
\end{definition}

By Lemmas~\ref{tw:base}(d) and~\ref{tw:kummer-field}, $g\mapsto\bar g$
identifies $G_S/D_2G_S=\Gal(E/B)$ with $\Ftwo^8$, and $a\mapsto\kappa_a$
identifies $V$ with the dual of $\Ftwo^8$, with $\kappa_{\alpha_i}(g)=\bar g_i$.
We write vectors in $\Ftwo^8$ as strings $c_0c_1\cdots c_7$ of their
coordinates.

For a place $\nu$ of $B$, a choice of a place of $B_S$ above $\nu$ gives a
homomorphism $\mathcal G_\nu\to G_S$, the \emph{local map} at $\nu$; another
choice changes it by an inner automorphism of $G_S$. For a finite
place $\nu\notin S$, the local map kills inertia, and its image is generated
by a Frobenius element $\Frob_\nu$. Let $\Sigma=S\cup\{v_1,v_2\}$.

\begin{definition}\label{tw:local-generators}
Fix a place of $B_S$ above each place of $B$, and use it to define the local
maps. For $\nu\in\Sigma$ we fix generators of $\mathcal G_\nu$ as follows,
and use the same names for their images in $G_S$: $\iota_j$ at $v_j$
(Lemma~\ref{tw:local-groups}(a)), so that $\iota_1,\iota_2$ are complex
conjugations; $\tau_\nu,\varphi_\nu$ at
$\nu\in\{\mathfrak q_1,\mathfrak q_2,\mathfrak r_1,\mathfrak r_2\}$, with the
generator $\alpha_i$ of $\nu$ as uniformizer $\varpi$ in
Lemma~\ref{tw:local-groups}(b); and $x_j,y_j,z_j$ at $\mathfrak p_j$, the
images of one fixed triple $x,y,z\in\mathcal G_{\Q_2}=\mathcal G_{\mathfrak p_j}$
lifting the elements $x,y,z$ of Lemma~\ref{tw:local-two}(b). We also write
$\Frob_{\mathfrak t_k}$ for $k=1,2,3$. These elements are the \emph{local
generators}: those at $\nu\in\Sigma$ are the generators just fixed, and the
local generator at $\mathfrak t_k$, $k=1,2,3$, is $\Frob_{\mathfrak t_k}$. We
call $\mathfrak t_1,\mathfrak t_2,\mathfrak t_3$ the \emph{capped primes}.
\end{definition}

The triple $x,y,z$ satisfies the hypothesis of
Lemma~\ref{tw:local-groups}(c). For a finite place $\nu\notin\Sigma$, in
particular for $\nu=\mathfrak t_4$, the elementary image of $\Frob_\nu$ does
not depend on the place of $B_S$ chosen above $\nu$, because $\Gal(E/B)$ is
abelian; it is the elementary image of the Frobenius of $\nu$ in $E/B$.

\begin{table}[htbp]
\centering
\caption{Elementary images in $\Ftwo^8$ of the local generators. Entry $i$
is $1$ exactly when the element negates $\sqrt{\alpha_i}$. The third and
sixth columns give the action on $\sqrt b$ for $b$ in the completion.}
\label{tw:vector-table}
\begin{tabular}{lll@{\qquad}lll}
\toprule
element & image & action & element & image & action\\
\midrule
$\iota_1$ & $10111010$ & sign at $v_1$ & $\iota_2$ & $11000101$ & sign at $v_2$\\
$x_1$ & $00100000$ & $(5,b)_{\mathfrak p_1}$ & $x_2$ & $00010000$ & $(5,b)_{\mathfrak p_2}$\\
$y_1$ & $11110010$ & $(-1,b)_{\mathfrak p_1}$ & $y_2$ & $10000001$ & $(-1,b)_{\mathfrak p_2}$\\
$z_1$ & $10000100$ & $(-2,b)_{\mathfrak p_1}$ & $z_2$ & $11111000$ & $(-2,b)_{\mathfrak p_2}$\\
$\tau_{\mathfrak q_1}$ & $00001000$ & $(-1,b)_{\mathfrak q_1}$ & $\varphi_{\mathfrak q_1}$ & $10100111$ & $(-\alpha_4,b)_{\mathfrak q_1}$\\
$\tau_{\mathfrak q_2}$ & $00000100$ & $(-1,b)_{\mathfrak q_2}$ & $\varphi_{\mathfrak q_2}$ & $11101011$ & $(-\alpha_5,b)_{\mathfrak q_2}$\\
$\tau_{\mathfrak r_1}$ & $00000010$ & $(2,b)_{\mathfrak r_1}$ & $\varphi_{\mathfrak r_1}$ & $01100101$ & $(-\alpha_6,b)_{\mathfrak r_1}$\\
$\tau_{\mathfrak r_2}$ & $00000001$ & $(2,b)_{\mathfrak r_2}$ & $\varphi_{\mathfrak r_2}$ & $01011010$ & $(-\alpha_7,b)_{\mathfrak r_2}$\\
$\Frob_{\mathfrak t_1}$ & $00100111$ & $(29,b)_{\mathfrak t_1}$ & $\Frob_{\mathfrak t_2}$ & $00011011$ & $(29,b)_{\mathfrak t_2}$\\
$\Frob_{\mathfrak t_3}$ & $00111000$ & $(41,b)_{\mathfrak t_3}$ & & &\\
\bottomrule
\end{tabular}
\end{table}

\begin{lemma}\label{tw:vectors}
The elementary images of the local generators are those of
Table~\ref{tw:vector-table}. More precisely, for each local generator $g$ at
a place $\nu$, the column ``action'' gives an element $c\in B_\nu^\times$
($c=-1$ at a real place, where $(-1,a)_\nu$ is the sign of $a$) such that
$g(\sqrt a)=(c,a)_\nu\sqrt a$ for every $a\in V$; hence
$\bar g=\bigl(\tfrac12(1-(c,\alpha_i)_\nu)\bigr)_{0\le i\le7}$. The
elementary image of the Frobenius $\Frob_{\mathfrak t_4}$ of the uncapped
prime $\mathfrak t_4$ above $41$ is $00110100$.
\end{lemma}

\begin{proof}
The actions in the column ``action'' follow from
Lemma~\ref{tw:local-groups} and \eqref{tw:hilbert}. At a real place,
$\iota_j$ negates $\sqrt a$ exactly when $a<0$ there, that is, when
$(-1,a)_{v_j}=-1$. At a place $\nu$ above $p\in\{3,5\}$, by
Lemma~\ref{tw:local-groups}(b), $\tau_\nu$ negates $\sqrt a$ exactly when
$\operatorname{ord}_\nu(a)$ is odd, that is, when $(u,a)_\nu=-1$ for the
nonsquare unit $u$ of Table~\ref{tw:squareclass-table}; and $\varphi_\nu$
fixes the square root of the generator $\alpha_i$ of $\nu$ and negates the
square root of a unit exactly when it is not a square modulo $\nu$, which is
the action of $(-\alpha_i,\cdot)_\nu$. At $\mathfrak p_j$ the actions are
those of Lemma~\ref{tw:local-groups}(c). At a capped prime $\nu$ above $p$,
every $a\in V$ is a unit, and the Frobenius negates $\sqrt a$ exactly when
$a$ is not a square modulo $\nu$, that is, when $(p,a)_\nu=-1$. These
Hilbert symbols are the local Artin symbols
\cite[Chapter~III, Remark~4.5]{MilneCFT}. The vector $\bar g$
is then read off from the square classes of the $\alpha_i$ in
Table~\ref{tw:squareclass-table} and \eqref{tw:hilbert}, and at the primes
$\mathfrak t_3,\mathfrak t_4$ above $41$ also from Lemma~\ref{tw:base}(e). For example, at the places above $3$ and
$5$ the vector of $\tau_\nu$ marks the odd valuations, and that of
$\varphi_\nu$ marks the units that are not squares modulo $\nu$, with a zero
at the generator of $\nu$. By Lemma~\ref{tw:base}(e), the $\alpha_i$ that are
not squares modulo $\mathfrak t_4$ are those with $i=2,3,5$, which gives
$00110100$ for $\Frob_{\mathfrak t_4}$. The supplementary program
\texttt{kummer241.gp} checks the norms of the $\alpha_i$, the generators of
the primes in Lemma~\ref{tw:base}(c) and the residues of $\sqrt{241}$ in
Table~\ref{tw:squareclass-table}, and recomputes
Table~\ref{tw:vector-table} from these Hilbert symbols, except the row of
$\mathfrak t_3$, which \texttt{kummer41.gp} computes in the same way; it also
computes the elementary image of $\Frob_{\mathfrak t_4}$.
\end{proof}

Fix a minimal presentation $1\to R\to\mathfrak F\to G_S\to1$, where
$\mathfrak F$ is a free pro-$2$ group of rank $8$. Minimality means
$R\subset D_2\mathfrak F$, so $\mathfrak F/D_2\mathfrak F=\Ftwo^8$ with the
coordinates above. Put $\mathfrak L=\operatorname{gr}\mathfrak F$, so that
$\dim\mathfrak L_1=8$, $\dim\mathfrak L_2=36$ and $\dim\mathfrak L_3=168$.
For $v,w\in\mathfrak L_1$ we have $(v+w)^{[2]}=v^{[2]}+w^{[2]}+[v,w]$, and the
square of an element of $\mathfrak F$ with image $v$ has class $v^{[2]}$. For
each local generator $g$ above choose a lift $\hat g\in\mathfrak F$, and define
the local relators
\begin{equation}\label{tw:relators}
 \rho_{v_j}=\hat\iota_j^{\,2},\qquad
 \rho_\nu=\hat\varphi_\nu\hat\tau_\nu\hat\varphi_\nu^{-1}\hat\tau_\nu^{-p}
 \ (\nu\in\{\mathfrak q_1,\mathfrak q_2,\mathfrak r_1,\mathfrak r_2\}),\qquad
 \rho_{\mathfrak p_j}=r(\hat x_j,\hat y_j,\hat z_j),
\end{equation}
where $p\in\{3,5\}$ is the rational prime below $\nu$, and $r$ is a
defining relator of $\mathcal G_{\Q_2}$ with respect to the fixed triple
$x,y,z$, the same for $j=1,2$. They lie in $R$, because each local relation
holds in $\mathcal G_\nu$. Their classes in $\mathfrak L_2$ are obtained by
substituting elementary images into the quadratic initials of
Lemma~\ref{tw:local-groups}:
\begin{equation}\label{tw:local-initials}
 \bar\iota_j^{\,[2]},\qquad
 [\bar\tau_\nu,\bar\varphi_\nu]+\tfrac{p-1}2\bar\tau_\nu^{\,[2]},\qquad
 \bar y_j^{\,[2]}+[\bar x_j,\bar y_j]+[\bar x_j,\bar z_j].
\end{equation}
The results below hold for every minimal presentation and every choice of
lifts.

The classes \eqref{tw:local-initials} have an arithmetic meaning, which we
now describe (Lemmas~\ref{tw:beta} and~\ref{tw:local-forms}). As in the
conventions, we realize $\mathfrak L$ inside the free associative algebra
$\Ftwo\langle X_0,\dots,X_7\rangle$, where $X_0,\dots,X_7$ is the standard
basis of $\mathfrak L_1=\Ftwo^8$. Through the identification of $\mathfrak L_1$
with $\Gal(E/B)$, the Kummer characters $\kappa_a$, $a\in V$, are linear forms
on $\mathfrak L_1$, with $\kappa_{\alpha_i}(X_j)=1$ exactly when $i=j$.

\begin{lemma}\label{tw:beta}
There is a unique linear isomorphism $\beta$ from $\mathfrak L_2$ onto the
space of symmetric bilinear forms on $V$ such that
\[
 \beta(v^{[2]})(a,b)=\kappa_a(v)\kappa_b(v),\qquad
 \beta([v,w])(a,b)=\kappa_a(v)\kappa_b(w)+\kappa_a(w)\kappa_b(v)
\]
for all $v,w\in\mathfrak L_1$ and $a,b\in V$. If $\xi\in\mathfrak L_2$ is
written in the free algebra as $\sum_{i,j}n_{ij}X_iX_j$, then
$\beta(\xi)(\alpha_i,\alpha_j)=n_{ij}$.
\end{lemma}

\begin{proof}
The elements $X_i^{[2]}=X_i^2$ and $[X_i,X_j]=X_iX_j+X_jX_i$ ($i<j$) form a
basis of $\mathfrak L_2$. Define $\beta$ on this basis by the displayed
formulas for basis vectors $v,w$; it sends $X_i^{[2]}$ to the form that is $1$
at $(\alpha_i,\alpha_i)$ and $0$ at the other pairs of basis vectors, and
$[X_i,X_j]$ to the form that is $1$ at $(\alpha_i,\alpha_j)$ and
$(\alpha_j,\alpha_i)$ and $0$ elsewhere. These $36$ forms are a basis of the
symmetric bilinear forms on $V$, so $\beta$ is an isomorphism, and it is the
only linear map satisfying the formulas on basis vectors. The formulas then
hold for all $v,w$, because the bracket is bilinear,
$(v+w)^{[2]}=v^{[2]}+w^{[2]}+[v,w]$, and squaring is the identity on $\Ftwo$.
The last statement holds for the basis elements $X_i^2$ and
$X_iX_j+X_jX_i$ by the description of their images, and both sides are
linear in $\xi$.
\end{proof}

\begin{definition}\label{tw:hilbert-form}
For $\nu\in\Sigma$, the \emph{local Hilbert form} $H_\nu$ is the symmetric
bilinear form on $V$ with $H_\nu(a,b)=1$ if $(a,b)_\nu=-1$ and
$H_\nu(a,b)=0$ otherwise.
\end{definition}

\begin{lemma}\label{tw:local-forms}
For every $\nu\in\Sigma$, $\beta$ maps the class of $\rho_\nu$ to $H_\nu$. The
eight classes \eqref{tw:local-initials}, one for each $\nu\in\Sigma$, sum to
zero, and the seven classes with $\nu\ne\mathfrak p_1$ are linearly
independent.
\end{lemma}

\begin{proof}
For a local generator $g$ at $\nu$, $\kappa_a(\bar g)$ is the value at $g$ of
the local Kummer character of the image of $a$ in $B_\nu$. Hence
$\beta$ (Lemma~\ref{tw:beta}) of the class of $\rho_\nu$ is the pullback to $V$ of a symmetric
bilinear form on $B_\nu^\times/B_\nu^{\times2}$, as is $H_\nu$, and it
suffices to compare the two on a basis of $B_\nu^\times/B_\nu^{\times2}$. At a
real place the only class is $-1$, and both forms take the value $1$ on
$(-1,-1)$. At $\nu$ above $p\in\{3,5\}$, take the basis $\varpi,u$ with
$\varpi$ the generator of $\nu$ and $u$ a nonsquare unit. By
Lemma~\ref{tw:local-groups}(b), $\kappa_\varpi(\tau)=\kappa_u(\varphi)=1$ and
$\kappa_\varpi(\varphi)=\kappa_u(\tau)=0$, so the form of
$[\tau,\varphi]+\frac{p-1}2\tau^{[2]}$ takes the values $\frac{p-1}2$, $1$,
$0$ on $(\varpi,\varpi)$, $(\varpi,u)$, $(u,u)$; by \eqref{tw:hilbert} these
are the indicators of $(\varpi,\varpi)_p=(-1)^{(p-1)/2}$, $(\varpi,u)_p=-1$
and $(u,u)_p=1$. At $\mathfrak p_j$ the characters $\kappa_2,\kappa_{-5},\kappa_5$
are dual to $x,y,z$, as in the proof of Lemma~\ref{tw:local-groups}(c). So
$y^{[2]}$ contributes the entry at $(-5,-5)$, $[x,y]$ the entries at $(2,-5)$
and $(-5,2)$, and $[x,z]$ those at $(2,5)$ and $(5,2)$: the form of
$y^{[2]}+[x,y]+[x,z]$ on the basis $2,-5,5$ is the matrix of the Hilbert
symbol displayed in that proof, from which this initial was obtained.

For $a,b\in V$ and a place $\nu'\notin\Sigma$, the place $\nu'$ is finite
and not above $2$, and $a,b$ are units at $\nu'$, so $(a,b)_{\nu'}=1$
\cite[Chapter~VIII, 5.8]{MilneCFT}. By the
product formula for the Hilbert symbol \cite[Chapter~V, Example~5.4]{MilneCFT},
$\sum_{\nu\in\Sigma}H_\nu=0$, so the eight classes sum to zero.

Finally, by Table~\ref{tw:squareclass-table} and \eqref{tw:hilbert}, the
seven pairs
\[
 (-1,\alpha_2),\ (-1,\varepsilon),\ (-1,\alpha_7),\ (-1,\alpha_4),\
 (-1,\alpha_5),\ (\alpha_2,\alpha_6),\ (\varepsilon,\alpha_7)
\]
have Hilbert symbol $-1$ exactly at $\{\mathfrak p_1,v_1\}$,
$\{\mathfrak p_1,v_2\}$, $\{\mathfrak p_2,v_2\}$, $\{\mathfrak q_1,v_1\}$,
$\{\mathfrak q_2,v_2\}$, $\{\mathfrak r_1,v_1\}$ and $\{\mathfrak r_2,v_2\}$;
the supplementary program \texttt{cup241.gp} recomputes these symbols.
Evaluating a relation $\sum_{\nu\ne\mathfrak p_1}n_\nu H_\nu=0$, with
$n_\nu\in\Ftwo$, at the first two pairs gives $n_{v_1}=n_{v_2}=0$, and then
at the other five pairs gives
$n_{\mathfrak p_2}=n_{\mathfrak q_1}=n_{\mathfrak q_2}=n_{\mathfrak r_1}=n_{\mathfrak r_2}=0$.
\end{proof}

\begin{lemma}\label{tw:complete-global-presentation}
The group $G_S$ has generator rank $8$ and relation rank $7$. The
seven relators $\rho_\nu$ with $\nu\in\Sigma\setminus\{\mathfrak p_1\}$ form
a minimal system of defining relations: their normal closure in
$\mathfrak F$ is $R$. In particular the dyadic local relator
$\rho_{\mathfrak p_1}$ lies in the normal closure of the other seven.
\end{lemma}

\begin{proof}
The generator rank is $\dim H^1(G_S,\Ftwo)=8$, by
Lemma~\ref{tw:kummer-field} and
\cite[Proposition~3.9.1]{NeukirchSchmidtWingberg2008}.

Next we bound $H^2$. Let $X=\operatorname{Spec}\mathcal O_B[1/30]$. Since
$2$ is invertible on $X$, the sheaf $\mu_2$ is the constant sheaf $\Ftwo$.
The finite subextensions of $B_S/B$ correspond to the connected finite étale
Galois $2$-covers $Y\to X$, and the Hochschild--Serre spectral sequence of
this system of covers \cite[Theorem~14.9 and Remark~14.10]{MilneEtale} gives an exact
sequence
\[
 0\to H^1(G_S,\Ftwo)\to H^1(X,\Ftwo)\to
 \bigl(\varinjlim_Y H^1(Y,\Ftwo)\bigr)^{G_S}\to
 H^2(G_S,\Ftwo)\to H^2(X,\Ftwo).
\]
A class in $H^1(Y,\Ftwo)$ is an étale double cover of $Y$. Its Galois closure
over $X$ is a finite étale Galois cover whose group embeds in the wreath
product of $\Gal(Y/X)$ with $C_2$, hence a $2$-cover in this system, over which
the class vanishes. So the direct limit is zero, and
$H^2(G_S,\Ftwo)$ injects into $H^2(X,\mu_2)$. The Kummer sequence and
$\Pic(X)=0$ identify $H^2(X,\mu_2)$ with $\Br(X)[2]$, where
$\Br(X)=H^2(X,\mathbb G_m)$. By \cite[Theorem~6.8.3 and
Example~6.9.4]{Poonen}, $\Br(X)$ is the kernel of the sum of the local
invariants $\bigoplus_{\nu\in\Sigma}\Br(B_\nu)\to\Q/\Z$. Each of the eight
groups $\Br(B_\nu)[2]$ has order two, and the sum of the invariants maps their
direct sum onto $\frac12\Z/\Z$, so $\dim\Br(X)[2]=7$. We use only the
resulting bound $\dim H^2(G_S,\Ftwo)\le\dim\Br(X)[2]\le7$.

By Lemma~\ref{tw:local-forms}, the classes in $\mathfrak L_2$ of the seven
relators $\rho_\nu$, $\nu\ne\mathfrak p_1$, are linearly independent. Since
$R\subset D_2\mathfrak F$, we have $R^2[R,\mathfrak F]\subset D_3\mathfrak F$, so
the classes of these seven relators in $R/R^2[R,\mathfrak F]$ are independent.
On the other hand
$\dim R/R^2[R,\mathfrak F]=\dim H^2(G_S,\Ftwo)\le7$ by the five-term
exact sequence of the presentation
\cite[Proposition~3.9.5]{NeukirchSchmidtWingberg2008},
\cite[(3.3)--(3.4)]{Quadrelli2021}. Hence the seven classes form a basis, the
relation rank is $7$, and by the pro-$2$ Nakayama lemma
\cite[Corollary~3.9.3]{NeukirchSchmidtWingberg2008} the seven relators
generate $R$ as a normal subgroup. Finally $\rho_{\mathfrak p_1}\in R$.
\end{proof}

\begin{corollary}\label{tw:any-seven}
Any seven of the eight local relators $\rho_\nu$, $\nu\in\Sigma$, generate
$R$ as a normal subgroup.
\end{corollary}

\begin{proof}
By Lemma~\ref{tw:local-forms}, the only linear relation among the eight
classes \eqref{tw:local-initials} is that their sum vanishes, and this
relation is Hilbert reciprocity. Hence any seven of the eight classes are
linearly independent, and the proof of
Lemma~\ref{tw:complete-global-presentation} applies to any seven of the eight
local relators.
\end{proof}

We omit the dyadic relator at $\mathfrak p_1$.

\begin{remark}\label{tw:cup-remark}
By Lemma~\ref{tw:complete-global-presentation}, the bound
$\dim H^2(G_S,\Ftwo)\le\dim\Br(X)[2]=7$ is attained, so
$H^2(G_S,\Ftwo)\to\Br(X)[2]$ is an isomorphism. This can also be seen
directly: the cup product
$\kappa_a\cup\kappa_b$ maps to the class of the quaternion algebra $(a,b)$, whose
invariant at $\nu$ is given by $(a,b)_\nu$ \cite[Chapter~III,
Remark~4.7]{MilneCFT}, and the seven pairs in the proof of
Lemma~\ref{tw:local-forms} give seven classes with independent invariant
vectors. The supplementary program \texttt{cup241.gp} computes the invariant
vectors at the eight places of $\Sigma$ of all $36$ algebras
$(\alpha_i,\alpha_j)$.
\end{remark}

\subsection{The Prescribed Local Groups}\label{tw:prescribed-section}
This subsection defines the quotient $G_B$ of $G_S$ with prescribed local
groups and the tower $B_G/B$ (Definition~\ref{tw:GB}), and gives a
presentation of $G_B$ by eight generators and thirty relators
(Lemma~\ref{tw:GB-presentation}), from which
Subsections~\ref{tw:retained-section}--\ref{tw:infinite-section} compute.

We now cut $G_S$ down to prescribed local groups. At the dyadic
places we use the restriction
$\mathcal G_{\mathfrak p_j}=\mathcal G_{\Q_2}\to\Gal(L_2/\Q_2)=\mathcal D$ of
Lemma~\ref{tw:local-two}. At a place $\nu$ above $p\in\{3,5\}$, the maximal
elementary abelian quotient of $\mathcal G_\nu$ is
$\Gal(\Q_p(\sqrt u,\sqrt p)/\Q_p)\cong C_2\times C_2$ for a nonsquare unit $u$;
this quotient has ramification index and residue degree two. At
the capped primes $\mathfrak t_1,\mathfrak t_2$ and $\mathfrak t_3$, whose
completions are $\Q_{29}$, $\Q_{29}$ and $\Q_{41}$, we bound the order of the
Frobenius by four. (A variant of the construction caps the inert prime
$\mathfrak t_0$ instead of $\mathfrak t_3$; see
Remark~\ref{tw:same-quotient}.)

\begin{definition}\label{tw:GB}
Let $N$ be the normal closure in $G_S$ of the images of
$\ker(\mathcal G_{\mathfrak p_j}\to \mathcal D)$ for $j=1,2$, of the Frattini subgroups
$D_2\mathcal G_\nu$ for $\nu\in\{\mathfrak q_1,\mathfrak q_2,\mathfrak r_1,\mathfrak r_2\}$,
and of $\Frob_{\mathfrak t_k}^4$ for $k=1,2,3$. Put $G_B=G_S/N$ and
$\overline G_B=G_B/D_4G_B$, and let $B_G\subset B_S$ be the fixed field of
$N$, so that $\Gal(B_G/B)=G_B$. We call the extension $B_G/B$ the
\emph{tower}.
\end{definition}

This does not depend on the choice of places of $B_S$ above each $\nu$:
another choice changes each local map by an inner automorphism of $G_S$,
which replaces the images above by conjugates and does not change their
normal closure $N$.

\begin{lemma}\label{tw:GB-basic}
\begin{enumerate}
\item[(a)] The image of $\mathcal G_{\mathfrak p_j}$ in $G_B$ is a quotient of
 $\mathcal D$, the image of $\mathcal G_\nu$ for $\nu$ above $3$ or $5$ is a quotient of
 $C_2\times C_2$ in which inertia maps to the image of $\tau_\nu$, and
 $\Frob_{\mathfrak t_k}$ has order dividing four in $G_B$ for $k=1,2,3$.
\item[(b)] $N\subseteq D_2G_S$. Hence $G_B/D_2G_B=\Gal(E/B)$, and $E\subseteq B_G$
 is the fixed field of $D_2G_B$.
\end{enumerate}
\end{lemma}

\begin{proof}
(a) By Definition~\ref{tw:GB}, the maps $\mathcal G_{\mathfrak p_j}\to G_B$
and $\mathcal G_\nu\to G_B$ kill $\ker(\mathcal G_{\mathfrak p_j}\to \mathcal D)$ and
$D_2\mathcal G_\nu$, and $\Frob_{\mathfrak t_k}^4\in N$. For $\nu$ above $3$
or $5$, $\mathcal G_\nu/D_2\mathcal G_\nu\cong C_2\times C_2$, and the inertia
group of $\mathcal G_\nu$ is generated by $\tau_\nu$
(Lemma~\ref{tw:local-groups}(b)).

(b) Since $\mathcal D$ and $\mathcal G_{\Q_2}$ both have generator rank three,
$\ker(\mathcal G_{\mathfrak p_j}\to \mathcal D)$ lies in the Frattini subgroup of
$\mathcal G_{\mathfrak p_j}$. The images of the Frattini subgroups of the
local groups, and the fourth powers $\Frob_{\mathfrak t_k}^4$, lie in
$D_2G_S$. Hence $N\subseteq D_2G_S$, and $G_B/D_2G_B=G_S/D_2G_S=\Gal(E/B)$ by
Lemma~\ref{tw:kummer-field}.
\end{proof}

By Lemma~\ref{tw:retained-quotient}(b) below, the quotients in
Lemma~\ref{tw:GB-basic}(a) are $\mathcal D$ and $C_2\times C_2$, and
$\Frob_{\mathfrak t_k}$ has order exactly four in $G_B$.

\begin{lemma}\label{tw:GB-presentation}
The kernel $N_B$ of $\mathfrak F\to G_S\to G_B$ is the normal closure
of the following thirty elements:
\begin{enumerate}
\item[(i)] the seven relators $\rho_\nu$, $\nu\in\Sigma\setminus\{\mathfrak p_1\}$;
\item[(ii)] for $j=1,2$: $\hat x_j^{\,2}$, $[\hat x_j,\hat y_j]$,
 $[\hat x_j,\hat z_j]$, $[[\hat y_j,\hat z_j],\hat z_j]$, $\hat z_j^{\,4}$,
 $[\hat y_j,\hat z_j]^2$;
\item[(iii)] for $\nu\in\{\mathfrak q_1,\mathfrak q_2,\mathfrak r_1,\mathfrak r_2\}$:
 $\hat\tau_\nu^{\,2}$ and $\hat\varphi_\nu^{\,2}$;
\item[(iv)] for $k=1,2,3$: $\widehat{\Frob}_{\mathfrak t_k}^{\,4}$, where $\widehat{\Frob}_{\mathfrak t_k}$ is the chosen
 lift of $\Frob_{\mathfrak t_k}$.
\end{enumerate}
Moreover $N_B\subset D_2\mathfrak F$, and $\rho_{\mathfrak p_1}\in N_B$.
\end{lemma}

\begin{proof}
By Lemma~\ref{tw:complete-global-presentation}, it suffices to show that the
images in $G_S$ of the elements (ii)--(iv) generate $N$ as a normal
subgroup. Let $\mathfrak F_3$ be free on $x,y,z$, and let $N_{\mathcal D}$ be the normal
closure of the seven relators of $\mathcal D$ in Lemma~\ref{tw:local-two}, so that
$\mathfrak F_3/N_{\mathcal D}=\mathcal D$. The relator $r$ maps to $1$ in $\mathcal D$, so $r\in N_{\mathcal D}$, and
$\ker(\mathcal G_{\Q_2}\to \mathcal D)$ is the image of $N_{\mathcal D}$. Since
$N_{\mathcal D}\subset D_2\mathfrak F_3$, we have
$N_{\mathcal D}^2[N_{\mathcal D},\mathfrak F_3]\subset D_3\mathfrak F_3$. The seven relators of $\mathcal D$
span $N_{\mathcal D}/N_{\mathcal D}^2[N_{\mathcal D},\mathfrak F_3]$
\cite[Corollary~3.9.3]{NeukirchSchmidtWingberg2008}; write the class of $r$ as
a combination of them. The classes in $\operatorname{gr}_2\mathfrak F_3$ of
$x^2,y^2,[x,y],[x,z]$ are independent, the relators $z^4$, $[y,z]^2$ and
$[[y,z],z]$ lie in $D_3\mathfrak F_3$, and $r$ has class
$y^{[2]}+[x,y]+[x,z]$; so $y^2$ has coefficient one. Replacing $y^2$ by $r$
therefore still gives a spanning set, and by
\cite[Corollary~3.9.3]{NeukirchSchmidtWingberg2008} the elements $r$, $x^2$,
$[x,y]$, $[x,z]$, $[[y,z],z]$, $z^4$, $[y,z]^2$ generate $N_{\mathcal D}$ as a normal
subgroup. Since $r=1$ in $\mathcal G_{\Q_2}$, the other six generate
$\ker(\mathcal G_{\Q_2}\to \mathcal D)$. For $\nu$ above $p\in\{3,5\}$, the quotient
of $\mathcal G_\nu$ by the normal closure of $\tau^2,\varphi^2$ is generated
by two involutions which commute, because
$\varphi\tau\varphi^{-1}=\tau^p=\tau$ there; so that normal closure is
$D_2\mathcal G_\nu$. Finally $N\subseteq D_2G_S$
(Lemma~\ref{tw:GB-basic}(b)) and
$R\subset D_2\mathfrak F$ give $N_B\subset D_2\mathfrak F$, and
$\rho_{\mathfrak p_1}\in R\subset N_B$.
\end{proof}

\subsection{\texorpdfstring{The Quotient $G_B/D_4G_B$}{The Quotient G\_B/D\_4G\_B}}\label{tw:retained-section}
This subsection computes $\operatorname{gr}_1$, $\operatorname{gr}_2$ and
$\operatorname{gr}_3$ of $G_B$ (of dimensions $8$, $15$ and $26$), shows that
the prescribed local groups inject into $G_B/D_3G_B$, and counts the
conjugates of $\iota_1$ in $G_B/D_4G_B$
(Lemmas~\ref{tw:graded}--\ref{tw:retained-quotient}). These results give
the hypotheses of Lemma~\ref{tw:filtered-local-freeness} in the proof of
Proposition~\ref{tw:infinite}, the bound $\theta\ge\theta_*$ of
Theorem~\ref{tw:field-family}, and the space $\operatorname{gr}_2G_B$ used in
Sections~\ref{an:section} and~\ref{dh:section}. Throughout, $\mathfrak F$,
$R$ and $\mathfrak L=\operatorname{gr}\mathfrak F$ are as in
Subsection~\ref{tw:global-section}, and $N_B$ is the kernel of
$\mathfrak F\to G_B$ (Lemma~\ref{tw:GB-presentation}).

\begin{definition}\label{tw:relation-space}
For $n\ge1$, the \emph{relation space} $R_n\subseteq\mathfrak L_n$ of $G_B$
in degree $n$ is the image of $N_B\cap D_n\mathfrak F$ in
$\mathfrak L_n=D_n\mathfrak F/D_{n+1}\mathfrak F$.
\end{definition}

\begin{lemma}\label{tw:graded}
For every $n\ge1$, $D_nG_B$ is the image of $D_n\mathfrak F$, and the induced
map $\mathfrak L_n\to\operatorname{gr}_nG_B$ is surjective with kernel $R_n$;
these maps respect brackets and restricted squares. In particular
$\operatorname{gr}_1G_B=\Ftwo^8$ and
$\operatorname{gr}_nG_B\cong\mathfrak L_n/R_n$.
\end{lemma}

\begin{proof}
The map $\mathfrak F\to G_B$ is surjective, and a surjection maps $D_n$ onto
$D_n$, so $D_nG_B$ is the image of $D_n\mathfrak F$. The kernel of
$D_n\mathfrak F\to D_nG_B$ is $N_B\cap D_n\mathfrak F$, and the preimage of
$D_{n+1}G_B$ in $D_n\mathfrak F$ is
$(N_BD_{n+1}\mathfrak F)\cap D_n\mathfrak F=(N_B\cap D_n\mathfrak F)\,D_{n+1}\mathfrak F$.
Hence $\operatorname{gr}_nG_B=D_nG_B/D_{n+1}G_B$ is the quotient of
$\mathfrak L_n$ by the image of $N_B\cap D_n\mathfrak F$, which is $R_n$.
Brackets and restricted squares are induced by commutators and squares, which
the homomorphism $\mathfrak F\to G_B$ preserves. Finally
$N_B\subset D_2\mathfrak F$ (Lemma~\ref{tw:GB-presentation}), so $R_1=0$ and
$\operatorname{gr}_1G_B=\mathfrak L_1=\Ftwo^8$.
\end{proof}

Consider the following $22$ elements of $\mathfrak L_2$, written with the
vectors of Table~\ref{tw:vector-table}:
\begin{equation}\label{tw:initials-list}
\begin{aligned}
 &\bar\iota_1^{\,[2]},\ \bar\iota_2^{\,[2]};\qquad
 [\bar\tau_\nu,\bar\varphi_\nu]+\bar\tau_\nu^{\,[2]}\ (\nu=\mathfrak q_1,\mathfrak q_2),\quad
 [\bar\tau_\nu,\bar\varphi_\nu]\ (\nu=\mathfrak r_1,\mathfrak r_2);\\
 &\bar y_j^{\,[2]}+[\bar x_j,\bar y_j]+[\bar x_j,\bar z_j],\quad
 \bar x_j^{\,[2]},\quad [\bar x_j,\bar y_j],\quad [\bar x_j,\bar z_j]\quad(j=1,2);\\
 &\bar\tau_\nu^{\,[2]},\ \bar\varphi_\nu^{\,[2]}\quad
 (\nu=\mathfrak q_1,\mathfrak q_2,\mathfrak r_1,\mathfrak r_2).
\end{aligned}
\end{equation}
They are the classes of the eight local relators \eqref{tw:local-initials}
and of the fourteen quadratic elements $\hat x_j^{\,2}$,
$[\hat x_j,\hat y_j]$, $[\hat x_j,\hat z_j]$, $\hat\tau_\nu^{\,2}$,
$\hat\varphi_\nu^{\,2}$ of Lemma~\ref{tw:GB-presentation}.

\begin{lemma}\label{tw:quadratic-layer}
The space $R_2$ is spanned by the $22$ elements \eqref{tw:initials-list}, and
$\dim R_2=21$; their only linear relation is that the eight classes of local
relators sum to zero. Hence $\dim\operatorname{gr}_2G_B=15$. Consequently, if
$g\in G_B$ has nonzero elementary image $v$, then $g$ has order at least four
in $G_B/D_3G_B$ exactly when $v^{[2]}\notin R_2$; otherwise $g$ has order two
there.
\end{lemma}

\begin{proof}
Conjugates of an element of $D_2\mathfrak F$ have the same class modulo
$D_3\mathfrak F$, classes of products add, and $\mathfrak L_2$ is finite.
Hence, by Lemma~\ref{tw:GB-presentation}, $R_2$ is spanned by the classes of
the thirty generators of $N_B$. Those in (i) have the classes
\eqref{tw:local-initials}, and the quadratic elements of (ii) and (iii) have
the classes $\bar x_j^{\,[2]}$, $[\bar x_j,\bar y_j]$, $[\bar x_j,\bar z_j]$,
$\bar\tau_\nu^{\,[2]}$, $\bar\varphi_\nu^{\,[2]}$. The remaining generators,
$[[\hat y_j,\hat z_j],\hat z_j]$, $\hat z_j^{\,4}$, $[\hat y_j,\hat z_j]^2$ and
the fourth powers in (iv), lie in $D_3\mathfrak F$. The class of
$\rho_{\mathfrak p_1}$ also lies in $R_2$, since $\rho_{\mathfrak p_1}\in N_B$.
The supplementary program \texttt{lie241.py} computes the rank of the $22$
elements in the $36$-dimensional space $\mathfrak L_2$: it is $21$. The sum of
the eight local classes vanishes by Lemma~\ref{tw:local-forms}, that is, by
Hilbert reciprocity, so this is the only relation; none of the fourteen
quadratic elements of (ii) and (iii) is involved in it. By
Lemma~\ref{tw:graded}, $\dim\operatorname{gr}_2G_B=36-21=15$. For the last
statement, $g^2\in D_2G_B$ has class $v^{[2]}$ in
$\operatorname{gr}_2G_B=\mathfrak L_2/R_2$ (Lemma~\ref{tw:graded}), so
$g^2\in D_3G_B$ exactly when $v^{[2]}\in R_2$.
\end{proof}

\begin{lemma}\label{tw:retained-quotient}
\begin{enumerate}
\item[(a)] $\dim\operatorname{gr}_3G_B=26$, so $|\overline G_B|=2^{8+15+26}=2^{49}$.
\item[(b)] For $j=1,2$, the local map induces an injection
 $\mathcal D\to G_B/D_3G_B$, and $\operatorname{gr}_n\mathcal D\to\operatorname{gr}_nG_B$ is
 injective for $n=1,2$. For $\nu$ above $3$ or $5$, the image of
 $\mathcal G_\nu$ in $G_B$ is
 $\langle\tau_\nu\rangle\times\langle\varphi_\nu\rangle\cong C_2\times C_2$,
 and it injects into $G_B/D_2G_B$. For $k=1,2,3$, $\Frob_{\mathfrak t_k}$ has
 order exactly four in $G_B$ and in $G_B/D_3G_B$. For $j=1,2$, $\iota_j$ has
 order two in $G_B/D_2G_B$.
\item[(c)] The conjugacy class of $\iota_1$ in $\overline G_B$ has $2^{15}$
 elements.
\item[(d)] We have $\bar\iota_1\ne\bar\iota_2$. For each
 $\nu\in\{\mathfrak p_1,\mathfrak p_2,\mathfrak q_1,\mathfrak q_2,\mathfrak r_1,\mathfrak r_2,\mathfrak t_1,\mathfrak t_2,\mathfrak t_3\}$,
 neither $\bar\iota_1$ nor $\bar\iota_2$ lies in the span of the elementary
 images of the local generators at $\nu$.
\end{enumerate}
\end{lemma}

\begin{proof}
(a) Work in $\mathfrak F/D_4\mathfrak F$, in which $D_2\mathfrak F/D_4\mathfrak F$
is elementary abelian and $D_3\mathfrak F/D_4\mathfrak F$ is central. Let
$\rho_1,\dots,\rho_{21}$ be the generators of $N_B$ listed in (i), the first
three elements of (ii) for each $j$, and (iii) of
Lemma~\ref{tw:GB-presentation}. They lie in $D_2\mathfrak F$, and their classes
in $\mathfrak L_2$ are the elements of \eqref{tw:initials-list} other than the
class of $\rho_{\mathfrak p_1}$, which are independent by
Lemma~\ref{tw:quadratic-layer}. The generators
$\omega_j=[[\hat y_j,\hat z_j],\hat z_j]$ lie in $D_3\mathfrak F$, and the
remaining generators of $N_B$ lie in $D_4\mathfrak F$. For $\rho\in D_2\mathfrak F$ and $g\in\mathfrak F$ we have
$g^{-1}\rho g=\rho[\rho,g]$, and modulo $D_4\mathfrak F$ the commutator
$[\rho,g]$ is central and depends only, and linearly, on the classes of
$\rho$ in $\mathfrak L_2$ and of $g$ in $\mathfrak L_1$. Hence the image of
$N_B$ in $\mathfrak F/D_4\mathfrak F$ is the subgroup generated by the $\rho_i$,
the $[\rho_i,f]$ for $f$ in a basis of $\mathfrak F$, and the $\omega_j$. An
element $\prod_i\rho_i^{a_i}$ times an element of $D_3\mathfrak F$ lies in
$D_3\mathfrak F$ only if all $a_i=0$. Therefore $R_3$ is spanned by the
brackets $[\bar\rho_i,\xi]$, for $\xi$ in a basis of $\mathfrak L_1$, and by
$[[\bar y_j,\bar z_j],\bar z_j]$. The supplementary program
\texttt{lie241c.py} computes these elements in the free associative algebra,
where $[a,\xi]=a\xi+\xi a$: they span a space of dimension $142$ in the
$168$-dimensional space $\mathfrak L_3$. So $\dim\operatorname{gr}_3G_B=26$
by Lemma~\ref{tw:graded}, and $|\overline G_B|=2^{8+15+26}$ by
Lemmas~\ref{tw:graded} and~\ref{tw:quadratic-layer}.

(b) The local map $\mathcal G_{\mathfrak p_j}\to G_B$ factors through $\mathcal D$ by
Definition~\ref{tw:GB}, and the induced map $\mathcal D\to G_B$ preserves the
Zassenhaus filtrations. It is injective on $\operatorname{gr}_1$ because
$\bar x_j,\bar y_j,\bar z_j$ are independent, and on $\operatorname{gr}_2$
because $\bar z_j^{\,[2]}$ and $[\bar y_j,\bar z_j]$, the images of
$\overline{z^2}$ and $\overline{[y,z]}$, are independent modulo $R_2$; the
program \texttt{lie241.py} checks both facts. If $1\neq g\in \mathcal D$, let $n\le2$
be maximal with $g\in D_n(\mathcal D)$; then the image of $g$ is nonzero in
$\operatorname{gr}_nG_B$, so it is not in $D_3G_B$. For $\nu$ above $3$ or
$5$, the vectors $\bar\tau_\nu$ and $\bar\varphi_\nu$ are independent, and the
image of $\mathcal G_\nu$ is a quotient of $C_2\times C_2$. For the Frobenius
elements, \texttt{lie241.py} checks that
$\overline{\Frob}_{\mathfrak t_k}^{\,[2]}\notin R_2$ for $k=1,2$, and
\texttt{check41.py} checks it for $k=3$ (and for the elementary image of
$\Frob_{\mathfrak t_4}$ in Lemma~\ref{tw:vectors}),
using the space $R_2$ of \texttt{lie241.py}; so
Lemma~\ref{tw:quadratic-layer} gives order at least four in $G_B/D_3G_B$,
while Definition~\ref{tw:GB} gives order at most four in $G_B$. Finally
$\bar\iota_j\neq0$.

(c) Write $\overline G=\overline G_B$ and $\iota=\iota_1$. The map
$\mathfrak L_1\to\operatorname{gr}_2\overline G$, $v\mapsto[\bar\iota,v]$, has
rank $7$ by \texttt{lie241.py}, so its kernel is the line through
$\bar\iota$. The map
$\operatorname{gr}_2\overline G\to\operatorname{gr}_3\overline G$,
$u\mapsto[u,\bar\iota]$, has rank $8$ by \texttt{lie241c.py}. Since
$D_2\overline G$ is abelian and $D_3\overline G$ is central in $\overline G$,
the map $h\mapsto[h,\iota]$ is a homomorphism $D_2\overline G\to D_3\overline G$;
it factors through $\operatorname{gr}_2\overline G$ and equals the second map.
Its kernel is $C_{\overline G}(\iota)\cap D_2\overline G$, where
$C_{\overline G}(\iota)=\{g\in\overline G:g\iota=\iota g\}$ is the
centralizer of $\iota$. The kernel therefore has order
$2^{15+26-8}=2^{33}$. If $g\in C_{\overline G}(\iota)$, then
$[\bar\iota,\bar g]=0$ in $\operatorname{gr}_2\overline G$, so
$\bar g\in\{0,\bar\iota\}$; conversely $\iota\in C_{\overline G}(\iota)$.
Hence $|C_{\overline G}(\iota)|=2^{34}$, and the class of $\iota$ has
$2^{49-34}=2^{15}$ elements.

(d) This is a finite check on Table~\ref{tw:vector-table}, carried out by
\texttt{lie241.py}; at $\mathfrak t_3$ the span is $\{0,00111000\}$, which
contains neither $10111010$ nor $11000101$ (\texttt{check41.py}).
\end{proof}

\subsection{Filtered Fox Calculus}\label{tw:fox-section}
This subsection and the next prove that $G_B$ is infinite
(Proposition~\ref{tw:infinite}) by the Golod--Shafarevich method. Suppose
that $G_B$ is finite, and let $\Lambda=\Ftwo[G_B]$. The rows
(Definition~\ref{tw:row}) of the relators of $G_B$ in a presentation with
eight generators (Lemma~\ref{tw:GB-presentation}) span a $\Lambda$-module
$\mathcal M$ whose Hilbert polynomial is $1-(1-8t)h_\Lambda$
(Lemma~\ref{tw:fox-rows}). The module $\mathcal M$ is also the sum of the
modules spanned by the rows of the relations of the eleven local groups
$\Gamma$, each isomorphic to $C_2$, $C_2\times C_2$, $C_4$ or $\mathcal D$; for $0<t<1$
the Hilbert polynomial of each of these modules is at most
$\psi_\Gamma(t)h_\Lambda(t)$, with $\psi_\Gamma$ as in
Definition~\ref{tw:psi} (Lemma~\ref{tw:filtered-local-freeness}). The row of
the dyadic relator $\rho_{\mathfrak p_1}$ of \eqref{tw:relators} lies both in
the module of $\mathfrak p_1$ and in the sum of the other ten, and the module
it generates has Hilbert polynomial $s_{\mathcal D}h_\Lambda$, with $s_{\mathcal D}$ as in
Lemma~\ref{tw:local-fox}. Comparing these bounds gives
$1\le h_\Lambda(t)P_B(t)$ for $0<t<1$, with $P_B$ as in \eqref{tw:PB}; this
fails at $t=34/117$ (Lemma~\ref{tw:PB-value}).

Let $\Delta$ be a finite $2$-group, $\Lambda=\Ftwo[\Delta]$, and
$\mathfrak I$ its augmentation ideal. We use the filtered spaces and Hilbert
polynomials of Definition~\ref{tw:hilbert-def}; in particular
$\operatorname{Fil}^m\Lambda=\mathfrak I^m$. In this subsection $n$ denotes
the rank of a free pro-$2$ group (Definition~\ref{tw:row}), and filtration
indices are written $m$. For an integer $c\ge1$, $\Lambda^c(-1)$ denotes
$\Lambda^c$ with $\operatorname{Fil}^0=\Lambda^c$ and
$\operatorname{Fil}^m=(\mathfrak I^{m-1})^c$ for $m\ge1$, so that
$h_{\Lambda^c(-1)}=ct\,h_\Lambda$. Subspaces carry the induced filtration
unless stated otherwise. A linear map $f\colon M\to M'$ of filtered spaces is
\emph{strict} if $f(\operatorname{Fil}^mM)=\operatorname{Fil}^mM'$ for all
$m$. For $0<t<1$,
\begin{equation}\label{tw:cumulative}
 h_M(t)=(1-t)\sum_{m\ge1}\dim(M/\operatorname{Fil}^mM)\,t^{m-1}.
\end{equation}

\begin{lemma}\label{tw:hilbert-properties}
Let $0<t<1$.
\begin{enumerate}
\item[(h1)] If $f\colon M\to M'$ is surjective with
 $f(\operatorname{Fil}^mM)\subset\operatorname{Fil}^mM'$ for all $m$, then
 $h_{M'}(t)\le h_M(t)$.
\item[(h2)] If $U_1\subset U_2$ are subspaces of the same filtered space,
 then $h_{U_1}(t)\le h_{U_2}(t)$.
\item[(h3)] If $f\colon M\to M'$ is surjective and strict, then
 $h_M=h_{\ker f}+h_{M'}$.
\item[(h4)] For subspaces $U_1,U_2$ of a filtered space,
 \begin{equation}\label{tw:filtered-sum}
  h_{U_1+U_2}(t)\le h_{U_1}(t)+h_{U_2}(t)-h_{U_1\cap U_2}(t).
 \end{equation}
\end{enumerate}
\end{lemma}

\begin{proof}
Statements (h1)--(h3) follow from the definitions and \eqref{tw:cumulative}.
For (h4), apply (h3) to $U_1\oplus U_2\to U_1+U_2$, $(u_1,u_2)\mapsto u_1+u_2$,
whose kernel is $U_1\cap U_2$, with the image filtration
$\operatorname{Fil}^mU_1+\operatorname{Fil}^mU_2$ on
$U_1+U_2$; this filtration is contained in the induced one, and (h1) applies
to the identity map.
\end{proof}

\begin{definition}\label{tw:row}
Let $\mathfrak F$ be a free pro-$2$ group with basis $f_1,\dots,f_n$ and
$\pi\colon\mathfrak F\to\Delta$ a surjection. Since $\Lambda^n\rtimes\Delta$,
with $\Delta$ acting by left multiplication, is a finite $2$-group, there is
a unique homomorphism $\mathfrak F\to\Lambda^n\rtimes\Delta$ sending $f_i$ to
$(e_i,\pi(f_i))$, where $e_1,\dots,e_n$ is the standard basis of $\Lambda^n$.
We write it as $w\mapsto(\partial w,\pi(w))$ and call
$\partial w=(\partial_1w,\dots,\partial_nw)\in\Lambda^n$ the \emph{row} of $w$.
\end{definition}

For $w$ in the discrete free group on the $f_i$, the row
$\partial w\in\Lambda^n$ is the image of the vector of free derivatives of $w$
\cite[Section~2]{Fox1953}.

\begin{lemma}\label{tw:row-properties}
\begin{enumerate}
\item[(F1)] $\partial(ww')=\partial w+\pi(w)\partial w'$.
\item[(F2)] $\sum_i\partial_iw\,(\pi(f_i)-1)=\pi(w)-1$.
\item[(F3)] If $\pi(\rho)=1$, then $\partial(u\rho u^{-1})=\pi(u)\partial\rho$;
 and $\partial$ restricts to a continuous homomorphism on $\ker\pi$.
\item[(F4)] If $\mathfrak F'$ is free on $f'_1,\dots,f'_{n'}$,
 $u_1,\dots,u_{n'}\in\mathfrak F$, and $\partial'$ denotes the rows for the map
 $\mathfrak F'\to\Delta$, $f'_l\mapsto\pi(u_l)$, then
 $\partial(W(u_1,\dots,u_{n'}))=\sum_l\partial'_lW\cdot\partial u_l$ for
 $W\in\mathfrak F'$.
\end{enumerate}
\end{lemma}

\begin{proof}
(F1) is the multiplication of the semidirect product. For (F2), the pairs
$(a,g)$ with $\sum a_i(\pi(f_i)-1)=g-1$ form a subgroup containing the
images of the $f_i$. (F3) follows from (F1). For (F4), both sides are the first
components of homomorphisms $\mathfrak F'\to\Lambda^n\rtimes\Delta$ that agree
on the basis.
\end{proof}

\begin{lemma}\label{tw:fox-rows}
Let $\mathcal M=\{a\in\Lambda^n:\sum_ia_i(\pi(f_i)-1)=0\}$, with the
filtration induced from $\Lambda^n(-1)$. Then
$h_{\mathcal M}=1-(1-nt)h_\Lambda$ and $\mathcal M=\partial(\ker\pi)$. If
$\rho$ lies in the normal closure of a set $Y\subset\ker\pi$, then
$\partial\rho$ lies in the $\Lambda$-span of $\partial Y$. In particular,
$\mathcal M$ is the $\Lambda$-span of the rows of any set of normal
generators of $\ker\pi$.
\end{lemma}

\begin{proof}
Since the $\pi(f_i)$ generate $\Delta$, we have
$\mathfrak I=\sum_i\Lambda(\pi(f_i)-1)\subset\sum_i\Ftwo(\pi(f_i)-1)+\mathfrak I^2$,
hence $\mathfrak I^m=\sum_i\mathfrak I^{m-1}(\pi(f_i)-1)$ for all $m\ge1$, by
induction and because $\mathfrak I$ is nilpotent. So
$\Lambda^n(-1)\to\mathfrak I$, $a\mapsto\sum a_i(\pi(f_i)-1)$, is strict and
surjective with kernel $\mathcal M$, and (h3) gives
$h_{\mathcal M}=nt\,h_\Lambda-(h_\Lambda-1)$.

By (F2), $\partial(\ker\pi)\subset\mathcal M$. Conversely, let
$\mathfrak F^\circ\subset\mathfrak F$ be the dense subgroup generated by the
$f_i$, which is a free group on them. Then $\pi(\mathfrak F^\circ)=\Delta$, and
the kernel of $\Ftwo[\mathfrak F^\circ]\to\Lambda$ is
$\Ftwo[\mathfrak F^\circ]\,(\mathfrak F^\circ\cap\ker\pi-1)$. Lift
$a\in\mathcal M$ to $\tilde a\in\Ftwo[\mathfrak F^\circ]^n$. Then
$\sum_i\tilde a_i(f_i-1)=\sum_l\zeta_l(q_l-1)$ with
$q_l\in\mathfrak F^\circ\cap\ker\pi$ and $\zeta_l\in\Ftwo[\mathfrak F^\circ]$.
By Fox's fundamental formula
$q_l-1=\sum_i(\partial q_l/\partial f_i)(f_i-1)$, where the
$\partial/\partial f_i$ are the free derivatives \cite[Section~2]{Fox1953},
and the coefficients of the $f_i-1$ are unique, because $\partial/\partial f_j$
maps $\sum_i\zeta'_i(f_i-1)$ to $\zeta'_j$. So
$\tilde a_i=\sum_l\zeta_l\,\partial q_l/\partial f_i$. On $\mathfrak F^\circ$
the rows are the images of the vectors of free derivatives, so
$a=\sum_l\bar\zeta_l\,\partial q_l$ lies in $\Lambda\,\partial(\ker\pi)$, which
equals $\partial(\ker\pi)$ by (F3). For the remaining statements, let $M_Y$ be
the $\Lambda$-span of $\partial Y$. By (F3), the elements $q\in\ker\pi$ with
$\partial q\in M_Y$ form a closed normal subgroup of $\mathfrak F$ containing
$Y$.
\end{proof}

The next lemma bounds, in terms of the function $\psi_\Gamma$ below, the
Hilbert polynomial of the $\Lambda$-module spanned by the rows of all
relations of one local group $\Gamma$. The groups
$\Gamma$ below are $C_2$, $C_2\times C_2$, $C_4$ and $\mathcal D$, with standard
generators $\iota$; $\tau,\varphi$; a generator; and $x,y,z$. Let $k$ be the
number of standard generators, so $k=1,2,1,3$. The Hilbert polynomials
$h_{\Ftwo[\Gamma]}$ are $1+t$, $(1+t)^2$, $(1+t)(1+t^2)$ and $h_{\Ftwo[\mathcal D]}$
(Theorem~\ref{tw:jennings}), and $D_3(\Gamma)=1$ in each case.

\begin{definition}\label{tw:psi}
For each of these groups $\Gamma$, put
\[
 \psi_\Gamma(t)=kt-1+h_{\Ftwo[\Gamma]}(t)^{-1}.
\]
\end{definition}

\begin{lemma}\label{tw:filtered-local-freeness}
Let $\Gamma\le\Delta$ be a subgroup isomorphic to one of these groups, with
$g_1,\dots,g_k\in\Gamma$ corresponding to its standard generators, and
suppose that the maps $\operatorname{gr}_1\Gamma\to\operatorname{gr}_1\Delta$
and $\operatorname{gr}_2\Gamma\to\operatorname{gr}_2\Delta$ are injective. Let $\Lambda_\Gamma=\Ftwo[\Gamma]\subset\Lambda$, with
augmentation ideal $\mathfrak I_\Gamma$, and let
$\mathcal M_\Gamma\subset\Lambda_\Gamma^k(-1)$ be the kernel of
$b\mapsto\sum_lb_l(g_l-1)$.
\begin{enumerate}
\item[(a)] There are elements $\eta_j\in\Lambda$ and integers $w_j\ge0$ such
 that $\Lambda=\bigoplus_j\eta_j\Lambda_\Gamma$ and
 $\mathfrak I^m=\bigoplus_j\eta_j\mathfrak I_\Gamma^{m-w_j}$ for every $m$,
 where $\mathfrak I_\Gamma^{m'}=\Lambda_\Gamma$ for $m'\le0$; moreover
 $\sum_jt^{w_j}=h_\Lambda/h_{\Ftwo[\Gamma]}$.
\item[(b)] The $\Lambda$-module $\Lambda\mathcal M_\Gamma\subset\Lambda^k(-1)$
 has $h_{\Lambda\mathcal M_\Gamma}=\psi_\Gamma h_\Lambda$.
\item[(c)] Let $\mathfrak F$ and $\pi$ be as above, let
 $u_1,\dots,u_k\in\mathfrak F$ with $\pi(u_l)=g_l$, and let $\mathfrak F_k$ be
 free on $k$ letters, mapping to $\Gamma$ by the standard generators. Let
 $U\subset\Lambda^n(-1)$ be the $\Lambda$-span of the rows
 $\partial W(u_1,\dots,u_k)$ for all $W$ in the kernel of
 $\mathfrak F_k\to\Gamma$. Then $h_U(t)\le\psi_\Gamma(t)h_\Lambda(t)$ for
 $0<t<1$.
\end{enumerate}
\end{lemma}

\begin{proof}
(a) Since $D_3(\Gamma)=1$ and $\operatorname{gr}\Gamma\to\operatorname{gr}\Delta$
is injective, $\Gamma\cap D_m\Delta=D_m(\Gamma)$ for all $m$. Choose elements
of $\Gamma$ whose classes form a homogeneous basis of $\operatorname{gr}\Gamma$,
and extend them by elements of $\Delta$ to a family whose classes form a
homogeneous basis of $\operatorname{gr}\Delta$. By Theorem~\ref{tw:jennings},
the ordered products of the elements $g-1$ of this family, each used at most
once and with the elements of $\Gamma$ last, form a basis of $\Lambda$, and
those of weight at least $m$ form a basis of $\mathfrak I^m$; here the weight
of a product is the sum of the degrees of the classes of its factors. The
same holds for $\Lambda_\Gamma$ and the elements of $\Gamma$ alone. Let the
$\eta_j$ be the ordered products of the other elements, and $w_j$ the weight
of $\eta_j$. The empty product $\eta_j=1$ has weight $0$, so
$\mathfrak I^m\cap\Lambda_\Gamma=\mathfrak I_\Gamma^m$: the filtration of
$\Lambda_\Gamma$ induced from $\Lambda$ is its augmentation filtration, with
Hilbert polynomial $h_{\Ftwo[\Gamma]}$. Comparing Hilbert polynomials gives the
last assertion.

(b) As in the proof of Lemma~\ref{tw:fox-rows}, the map
$\Lambda_\Gamma^k(-1)\to\mathfrak I_\Gamma$ is strict and surjective, so
$h_{\mathcal M_\Gamma}=kt\,h_{\Ftwo[\Gamma]}-(h_{\Ftwo[\Gamma]}-1)$. By (a),
$\Lambda^k(-1)=\bigoplus_j\eta_j\Lambda_\Gamma^k(-1)$ with
$\operatorname{Fil}^m=\bigoplus_j\eta_j\operatorname{Fil}^{m-w_j}$, and
$\Lambda\mathcal M_\Gamma=\bigoplus_j\eta_j\mathcal M_\Gamma$. Hence
$h_{\Lambda\mathcal M_\Gamma}=(h_\Lambda/h_{\Ftwo[\Gamma]})\,h_{\mathcal M_\Gamma}=\psi_\Gamma h_\Lambda$.

(c) For $W$ in the kernel, the local row $\partial'W\in\Lambda_\Gamma^k$ lies in
$\mathcal M_\Gamma$ by (F2), and by (F4)
$\partial W(u)=\sum_l\partial'_lW\,\partial u_l$. Thus $U$ is contained in the
image of $\Lambda\mathcal M_\Gamma$ under the map
$T_u\colon\Lambda^k(-1)\to\Lambda^n(-1)$, $a\mapsto\sum_la_l\partial u_l$,
which preserves filtrations because
$\partial u_l\in\Lambda^n=\operatorname{Fil}^0\Lambda^n$. By (h2),
(h1) and (b),
$h_U\le h_{T_u(\Lambda\mathcal M_\Gamma)}\le h_{\Lambda\mathcal M_\Gamma}=\psi_\Gamma h_\Lambda$.
\end{proof}

In the proof of Proposition~\ref{tw:infinite}, the row of the dyadic relator
$\rho_{\mathfrak p_1}$ lies both in the module spanned by the rows of the
relations at $\mathfrak p_1$ and in the module spanned by the rows of the
relations at the other ten places. The next lemma computes the Hilbert
polynomial of the module generated by this row.

\begin{lemma}\label{tw:local-fox}
In the situation of Lemma~\ref{tw:filtered-local-freeness} with $\Gamma=\mathcal D$,
suppose that $u_1,u_2,u_3$ are the basis elements $f_1,f_2,f_3$ of
$\mathfrak F$. Let $r\in\mathfrak F_3$ map to $1$ in $\mathcal D$, with class
$y^{[2]}+[x,y]+[x,z]$ in $\operatorname{gr}_2\mathfrak F_3$, and let
$\rho=r(f_1,f_2,f_3)$. Then
\[
 h_{\Lambda\partial\rho}=s_{\mathcal D}h_\Lambda,\qquad
 s_{\mathcal D}(t)=t^2\Bigl(1-\frac{t^7}{h_{\Ftwo[\mathcal D]}(t)}\Bigr).
\]
\end{lemma}

\begin{proof}
By (F4), $\partial\rho=(\partial'r,0,\dots,0)$, where
$\partial'r\in\Lambda_{\mathcal D}^3$ is the local row, so $\Lambda\partial\rho$ is the
module $\Lambda\partial'r\subset\Lambda^3(-1)$ placed in the first three
coordinates. By Lemma~\ref{tw:filtered-local-freeness}(a),
$\Lambda^3=\bigoplus_j\eta_j\Lambda_{\mathcal D}^3$ with
$\operatorname{Fil}^m(\Lambda^3(-1))=\bigoplus_j\eta_j\operatorname{Fil}^{m-w_j}(\Lambda_{\mathcal D}^3(-1))$,
and
$b\mapsto\eta_jb$ is injective on $\Lambda_{\mathcal D}^3$. The submodule
$\Lambda\partial'r=\bigoplus_j\eta_j\Lambda_{\mathcal D}\partial'r$ decomposes in the
same way, so
$\operatorname{Fil}^m(\Lambda\partial'r)=\bigoplus_j\eta_j\operatorname{Fil}^{m-w_j}(\Lambda_{\mathcal D}\partial'r)$
for
the induced filtrations, and
$h_{\Lambda\partial\rho}=(h_\Lambda/h_{\Ftwo[\mathcal D]})\,h_{\Lambda_{\mathcal D}\partial'r}$. It
therefore suffices to show that $h_{\Lambda_{\mathcal D}\partial'r}=t^2(h_{\Ftwo[\mathcal D]}-t^7)$, for
the filtration induced from $\Lambda_{\mathcal D}^3(-1)$. Write $\mathfrak I_{\mathcal D}$ for the augmentation ideal of
$\Lambda_{\mathcal D}=\Ftwo[\mathcal D]$, and $X_1,X_2,X_3,X_4$ for $x-1,y-1,z-1,w-1$, where
$w=[y,z]$.

\emph{The linear part of $\partial'r$.} Here $\pi$ also denotes the map
$\mathfrak F_3\to \mathcal D$. For $g,h\in\mathfrak F_3$, (F1) gives
$\partial'(g^2)=(1+\pi(g))\partial'g$ and
$\partial'[g,h]=\pi(g)^{-1}(\pi(h)^{-1}-1)\partial'g+\pi(g)^{-1}\pi(h)^{-1}(\pi(g)-1)\partial'h$.
Moreover $\partial'g$ is congruent modulo $\mathfrak I_{\mathcal D}$ to the image
$\bar g\in\Ftwo^3$ of $g$ in $\mathfrak F_3/D_2\mathfrak F_3$, and
$\pi(g)-1\equiv\sum_l\bar g_lX_l$ modulo $\mathfrak I_{\mathcal D}^2$ by (F2). Hence, for
$q\in D_2\mathfrak F_3$, the row $\partial'q$ has entries in $\mathfrak I_{\mathcal D}$,
and modulo $\mathfrak I_{\mathcal D}^2$ it is additive in $q$, because
$\pi(q)-1\in\mathfrak I_{\mathcal D}^2$. It vanishes modulo $\mathfrak I_{\mathcal D}^2$ on
$D_3\mathfrak F_3=(D_2\mathfrak F_3)^2[D_2\mathfrak F_3,\mathfrak F_3]$ (Lazard's
formula): for $q\in D_2\mathfrak F_3$ the formulas above give
$\partial'(q^2),\partial'[q,h]\in\mathfrak I_{\mathcal D}^2$, since $\bar q=0$ and
$\pi(q)-1\in\mathfrak I_{\mathcal D}^2$. So the row modulo $\mathfrak I_{\mathcal D}^2$ depends
only on the class of $q$ in $\operatorname{gr}_2\mathfrak F_3$, the classes
$v^{[2]}$ and $[v,v']$ giving the rows with $l$-th entries $v_l\sum_iv_iX_i$
and $v_l\sum_iv'_iX_i+v'_l\sum_iv_iX_i$.
For the class $y^{[2]}+[x,y]+[x,z]$ this gives
\[
 \partial'r\equiv(X_2+X_3,\ X_2+X_1,\ X_1)\pmod{\mathfrak I_{\mathcal D}^2}.
\]

\emph{Three facts about $\Lambda_{\mathcal D}$.} Let $A=\bigoplus_m\mathfrak I_{\mathcal D}^m/\mathfrak I_{\mathcal D}^{m+1}$,
and write $\bar X_1,\bar X_2,\bar X_3$ and $\bar X_4$ for the classes of $X_1,X_2,X_3$ in degree one and
of $X_4$ in degree two. By Theorem~\ref{tw:jennings} and
Lemma~\ref{tw:local-two}(c), the products
$X_1^{a_1}X_2^{a_2}X_3^{a_3}X_4^{a_4}$ with $a_1,a_2,a_4\in\{0,1\}$ and
$0\le a_3\le3$, of weight $a_1+a_2+a_3+2a_4\ge m$, form a basis of
$\mathfrak I_{\mathcal D}^m$; here $X_3^2=z^2-1$. Hence the monomials
$\bar X_1^{a_1}\bar X_2^{a_2}\bar X_3^{a_3}\bar X_4^{a_4}$ form a basis of $A$, and, since the
coefficients of $h_{\Ftwo[\mathcal D]}$ are $1,3,5,7,7,5,3,1$,
\[
 \dim\mathfrak I_{\mathcal D}^m=32,31,28,23,16,9,4,1,0\qquad(m=0,\dots,8).
\]
Because $x$ and $w$ are central and $x^2=y^2=w^2=z^4=1$, the elements $\bar X_1$
and $\bar X_4$ are central in $A$ and $\bar X_1^2=\bar X_2^2=\bar X_4^2=\bar X_3^4=0$. Moreover $\bar X_3\bar X_2=\bar X_2\bar X_3+\bar X_4$,
because $X_3X_2+X_2X_3=zy+yz=zy(w-1)\equiv X_4\pmod{\mathfrak I_{\mathcal D}^3}$; by
induction $\bar X_3^j\bar X_2=\bar X_2\bar X_3^j+j\bar X_3^{j-1}\bar X_4$ for $j\ge1$.

First, $\mathfrak I_{\mathcal D}^7=\Ftwo\Omega_{\mathcal D}$, where $\Omega_{\mathcal D}=\sum_{g\in \mathcal D}g$:
$\mathfrak I_{\mathcal D}^7$ is one-dimensional and $(g-1)\mathfrak I_{\mathcal D}^7\subset\mathfrak I_{\mathcal D}^8=0$,
so its nonzero element is invariant under left multiplication by $\mathcal D$, and
hence equal to $\Omega_{\mathcal D}$.

Second, for $m<7$ the map $A_m\to A_{m+1}^3$, $\xi\mapsto(\xi\bar X_1,\xi\bar X_2,\xi\bar X_3)$,
is injective. Let $\xi=\sum_\mu b_\mu\mu$ over the basis monomials
$\mu=\bar X_1^{a_1}\bar X_2^{a_2}\bar X_3^{a_3}\bar X_4^{a_4}$ of degree $m$. Right multiplication by $\bar X_1$
maps the monomials with $a_1=0$ to distinct basis monomials and kills the
others, and right multiplication by $\bar X_3$ maps those with $a_3\le2$ to
distinct basis monomials and kills those with $a_3=3$. So $\xi\bar X_1=\xi\bar X_3=0$
forces $\xi=\sum_{a_2,a_4}b_{a_2a_4}\bar X_1\bar X_2^{a_2}\bar X_3^3\bar X_4^{a_4}$. Then
\begin{align*}
 \xi\bar X_2={}&b_{00}(\bar X_1\bar X_2\bar X_3^3+\bar X_1\bar X_3^2\bar X_4)
 +b_{01}\bar X_1\bar X_2\bar X_3^3\bar X_4\\
 &+b_{10}\bar X_1\bar X_2\bar X_3^2\bar X_4,
\end{align*}
a combination of distinct basis monomials, so $\xi\bar X_2=0$ forces
$\xi=b_{11}\bar X_1\bar X_2\bar X_3^3\bar X_4$, which has degree $7$. Hence $\xi=0$ when $m<7$.

Third, the same holds for $\xi\mapsto(\xi(\bar X_2+\bar X_3),\xi(\bar X_2+\bar X_1),\xi\bar X_1)$, because this
triple is obtained from $(\bar X_1,\bar X_2,\bar X_3)$ by the invertible matrix
\[
 \begin{pmatrix}0&1&1\\1&1&0\\1&0&0\end{pmatrix}.
\]
This is the coefficient matrix of $y^{[2]}+[x,y]+[x,z]$, that is, the matrix
of the $2$-adic Hilbert symbol on $2,-5,5$ in the proof of
Lemma~\ref{tw:local-groups}(c). So the injectivity used here, and with it
the correction term $s_{\mathcal D}$, comes from the nondegeneracy of the local Hilbert pairing,
that is, from local duality. The supplementary program \texttt{dyadic241.py}
confirms these three facts by direct computation in $\Lambda_{\mathcal D}$.

\emph{The filtration.} If $\xi\in\mathfrak I_{\mathcal D}^m\setminus\mathfrak I_{\mathcal D}^{m+1}$
with $m<7$, then by the third fact $\xi\,\partial'r$ lies in $\operatorname{Fil}^{m+2}$ but not
in $\operatorname{Fil}^{m+3}$, whereas $\Omega_{\mathcal D}\partial'r=0$ because $\Omega_{\mathcal D}\mathfrak I_{\mathcal D}=0$.
Now let $\xi\in\Lambda_{\mathcal D}$ with $\xi\,\partial'r\in \operatorname{Fil}^{m+2}$, and suppose that
$\xi\notin\mathfrak I_{\mathcal D}^m+\Ftwo\Omega_{\mathcal D}$. Let $m'<m$ be maximal with
$\xi\in\mathfrak I_{\mathcal D}^{m'}+\Ftwo\Omega_{\mathcal D}$, and write $\xi=\xi'+a\,\Omega_{\mathcal D}$
with $\xi'\in\mathfrak I_{\mathcal D}^{m'}$ and $a\in\Ftwo$. Then
$\xi'\notin\mathfrak I_{\mathcal D}^{m'+1}$, and $m'<7$ because
$\mathfrak I_{\mathcal D}^7=\Ftwo\Omega_{\mathcal D}$; so $\xi\,\partial'r=\xi'\partial'r$ is not in
$\operatorname{Fil}^{m'+3}\supset \operatorname{Fil}^{m+2}$, a contradiction. Therefore
$\operatorname{Fil}^{m+2}(\Lambda_{\mathcal D}\partial'r)=\mathfrak I_{\mathcal D}^m\partial'r$ for every $m\ge0$,
the map $\xi\mapsto\xi\,\partial'r$ has kernel $\Ftwo\Omega_{\mathcal D}$, and
$h_{\Lambda_{\mathcal D}\partial'r}=\sum_{m<7}\dim(\mathfrak I_{\mathcal D}^m/\mathfrak I_{\mathcal D}^{m+1})\,t^{m+2}=t^2(h_{\Ftwo[\mathcal D]}-t^7)$.
\end{proof}

\subsection{Infinitude}\label{tw:infinite-section}
This subsection assembles the local bounds of
Subsection~\ref{tw:fox-section} into the Golod--Shafarevich function $P_B$
of \eqref{tw:PB}, shows that $P_B(34/117)<0$ (Lemma~\ref{tw:PB-value}), and
concludes that $G_B$ is infinite (Proposition~\ref{tw:infinite}). Theorem~\ref{tw:field-family}
uses this to obtain fields of every degree $2^n$ over $B$ with $n\ge49$.

By Definition~\ref{tw:psi} and the Hilbert polynomials listed before it,
\begin{equation}\label{tw:costs}
\begin{gathered}
 \psi_{C_2}(t)=\frac{t^2}{1+t},\qquad
 \psi_{C_2\times C_2}(t)=2t-1+\frac1{(1+t)^2},\\
 \psi_{C_4}(t)=\frac{t^4}{(1+t)(1+t^2)},\qquad
 \psi_{\mathcal D}(t)=3t-1+\frac1{h_{\Ftwo[\mathcal D]}(t)}.
\end{gathered}
\end{equation}
With $s_{\mathcal D}$ as in Lemma~\ref{tw:local-fox}, the \emph{Golod--Shafarevich
function} of $G_B$ is
\begin{equation}\label{tw:PB}
 P_B(t)=1-8t+2\psi_{C_2}(t)+4\psi_{C_2\times C_2}(t)+2\psi_{\mathcal D}(t)-s_{\mathcal D}(t)
 +3\psi_{C_4}(t).
\end{equation}
In the proof of Proposition~\ref{tw:infinite}, the term $-8t$ comes from the
eight generators of a free group mapping onto $G_B$; the terms in
$\psi_{C_2}$, $\psi_{C_2\times C_2}$, $\psi_{\mathcal D}$ and $\psi_{C_4}$ come from the
local relations at the two real places, the four places above $3$ and $5$,
the two dyadic places and the three places
$\mathfrak t_1,\mathfrak t_2,\mathfrak t_3$; and $-s_{\mathcal D}$ comes from the row of
the dyadic relator $\rho_{\mathfrak p_1}$, which lies both in the module of
$\mathfrak p_1$ and in the sum of the modules of the other places.

\begin{lemma}\label{tw:PB-value}
\[
 P_B\Bigl(\frac{34}{117}\Bigr)=-\frac{187433948535241}{88772225460489675}<0.
\]
\end{lemma}

\begin{proof}
By \eqref{tw:costs}, \eqref{tw:PB}, the formula for $s_{\mathcal D}$ in
Lemma~\ref{tw:local-fox} and $h_{\Ftwo[\mathcal D]}(t)=(1+t)^3(1+t^2)^2$
(Lemma~\ref{tw:local-two}(c)), $P_B$ is a rational function with rational
coefficients. Exact rational arithmetic, carried out by the supplementary
program \texttt{gs241.py}, gives the stated value.
\end{proof}

\begin{proposition}\label{tw:infinite}
The group $G_B$ is infinite.
\end{proposition}

\begin{proof}
Suppose that $G_B$ is finite, and put $\Delta=G_B$ and
$\Lambda=\Ftwo[\Delta]$. By Table~\ref{tw:vector-table},
$\bar x_1,\bar y_1,\bar z_1$ are independent; choose
$g_4,\dots,g_8\in G_S$ such that
$\bar x_1,\bar y_1,\bar z_1,\bar g_4,\dots,\bar g_8$ is a basis of
$\Ftwo^8$. Let $\mathfrak F$ be free on $f_1,\dots,f_8$, and send
$f_1,f_2,f_3$ to $x_1,y_1,z_1$ and $f_l$ to $g_l$ for $4\le l\le8$. This is a
minimal presentation of
$G_S$; we take the lifts $\hat x_1=f_1$, $\hat y_1=f_2$,
$\hat z_1=f_3$, and arbitrary lifts of the other local generators. Let
$\pi\colon\mathfrak F\to\Delta$ be the composite and $\mathcal M$ as in
Lemma~\ref{tw:fox-rows}, so that $h_{\mathcal M}=1-(1-8t)h_\Lambda$.

By Lemma~\ref{tw:retained-quotient}(b), the image in $\Delta$ of each
local group is the prescribed group $\Gamma$, and
$\operatorname{gr}_n\Gamma\to\operatorname{gr}_n\Delta$ is injective for
$n=1,2$; this is the hypothesis of
Lemma~\ref{tw:filtered-local-freeness}. Indeed, for $\Gamma=\mathcal D$ this is stated
in Lemma~\ref{tw:retained-quotient}(b). For $\Gamma=C_2$ and
$\Gamma=C_2\times C_2$ we have $\operatorname{gr}_2\Gamma=0$, and
$\operatorname{gr}_1\Gamma=\Gamma$ injects into $G_B/D_2G_B$ because
$\bar\iota_j\ne0$, respectively because the image of $\mathcal G_\nu$ injects
into $G_B/D_2G_B$. For $\Gamma=C_4$ generated by $\Frob_{\mathfrak t_k}$, $k=1,2,3$,
which has order four in $G_B/D_3G_B$, we have
$\Frob_{\mathfrak t_k}\notin D_2G_B$, since otherwise its square would lie in
$D_4G_B\subseteq D_3G_B$, and $\Frob_{\mathfrak t_k}^2\notin D_3G_B$; this is
the injectivity on $\operatorname{gr}_1\Gamma$ and on
$\operatorname{gr}_2\Gamma=\langle\Frob_{\mathfrak t_k}^2\rangle$. We use
eleven places: the real
places $v_1,v_2$, with $\Gamma=C_2$; the four places above $3$ and $5$, with
$\Gamma=C_2\times C_2$; the dyadic places $\mathfrak p_1,\mathfrak p_2$, with
$\Gamma=\mathcal D$; and $\mathfrak t_1,\mathfrak t_2,\mathfrak t_3$, with $\Gamma=C_4$
generated by the Frobenius. For each of these places $\nu$, let $U_\nu$ be
the module $U$ of Lemma~\ref{tw:filtered-local-freeness}(c) for this $\Gamma$,
formed with the lifts of the local generators at $\nu$. Every generator of
$N_B$ in Lemma~\ref{tw:GB-presentation} is of the form $W(u_1,\dots,u_k)$ for
one of these places and some $W$ in the kernel of $\mathfrak F_k\to\Gamma$, and
conversely every such element lies in $N_B$. By Lemma~\ref{tw:fox-rows},
$\mathcal M=\sum_\nu U_\nu$. Let $U'$ be the sum of the $U_\nu$ over the ten
places $\nu\ne\mathfrak p_1$. The row of $\rho_{\mathfrak p_1}$ lies in
$U_{\mathfrak p_1}$. Since $\rho_{\mathfrak p_1}$ lies in the normal closure of
the relators (i) of Lemma~\ref{tw:GB-presentation}, its row also lies in the
$\Lambda$-span of their rows (Lemma~\ref{tw:fox-rows}), which is contained in
$U'$. Hence $\Lambda\partial\rho_{\mathfrak p_1}\subset U_{\mathfrak p_1}\cap U'$.
By Lemma~\ref{tw:hilbert-properties}(h4) and (h2),
Lemma~\ref{tw:filtered-local-freeness}(c) and Lemma~\ref{tw:local-fox}, for
$0<t<1$,
\begin{align*}
 1-(1-8t)h_\Lambda(t)=h_{\mathcal M}(t)
 &\le h_{U_{\mathfrak p_1}}(t)+\sum_{\nu\ne\mathfrak p_1}h_{U_\nu}(t)
 -h_{\Lambda\partial\rho_{\mathfrak p_1}}(t)\\
 &\le\bigl(2\psi_{C_2}+4\psi_{C_2\times C_2}+2\psi_{\mathcal D}+3\psi_{C_4}-s_{\mathcal D}\bigr)(t)
 \,h_\Lambda(t),
\end{align*}
that is, $1\le h_\Lambda(t)P_B(t)$. Since $h_\Lambda(t)>0$, this contradicts
$P_B(34/117)<0$ (Lemma~\ref{tw:PB-value}).
\end{proof}

\subsection{The Fields}\label{tw:fields-section}
This subsection proves Theorem~\ref{tw:field-family}: the existence of the
fields $K$, their root discriminant $\rd(K)=\lambda$, the bound
$\theta\ge\theta_*$ on the signature of $F=K^{\langle\iota_1\rangle}$, and
the splitting and local types at the selected primes. These are the
properties of the fields used in
Sections~\ref{gf:section}--\ref{dh:section}, in Section~\ref{geo:section} and
in Section~\ref{sec:certificate}.

\begin{definition}\label{tw:rd}
The \emph{root discriminant} of a number field $M$ with discriminant
$\Delta_M$ is $\rd(M)=|\Delta_M|^{1/[M:\Q]}$.
\end{definition}

\begin{definition}\label{tw:type}
For a number field $M$ and a prime $\mathfrak P$ of $M$ above a rational
prime $p$, the \emph{absolute type} of $\mathfrak P$ is the pair $(e,f)$ of
its ramification index and residue degree over $p$. For a Galois extension
$M/A$ of number fields, all primes of $M$ above a prime $\mathfrak p$ of $A$
have the same ramification index $e$ and residue degree $f$ in $M/A$; the
pair $(e,f)$ is the \emph{type} of $\mathfrak p$ in $M/A$.
\end{definition}

As in the introduction, let $\lambda=2^{9/4}\sqrt{3615}$, where
$3615=3\cdot5\cdot241$.

\begin{theorem}\label{tw:field-family}
For every power of two $m\ge2^{49}$ there is a finite Galois extension
$K/B$ of degree $m$, contained in $B_G$, such that $G_B\to\overline G_B$
factors through $\Gal(K/B)$. In particular $K$ contains the fixed field of
$D_4G_B$, hence the fixed field of $D_3G_B$ and the Kummer field $E$. Every
such $K$ is totally imaginary. Let $\iota_1$ also denote complex
conjugation at any place of $K$ above $v_1$; these complex conjugations are
conjugate in $\Gal(K/B)$, one of them is the image of the element $\iota_1$
of Definition~\ref{tw:local-generators}, and the statements below do not
depend on the choice. Let $F=K^{\langle\iota_1\rangle}$ and
$d=[F:\Q]=m$, let $(b,c)$ be the signature of $F$ and $\theta=c/d$, and put
$\theta_*=\thetastar=1/2-2^{-17}$. Then $F$ has a real place,
\[
 \rd(K)=\rd(F)=\lambda,\qquad \theta_*\le\theta<\frac12 ,
\]
and $b/d=1/(2N_\iota)$, where $N_\iota\ge2^{15}$ is the number of conjugates
of $\iota_1$ in $\Gal(K/B)$. The extension $K/F$ is quadratic and unramified
at every finite place. Let $\mathfrak p\in S\cup\{\mathfrak t_1,\mathfrak
t_2,\mathfrak t_3\}$. Every prime of $F$ above $\mathfrak p$ splits in $K$, the
primes of $K$ and of $F$ above $\mathfrak p$ have the absolute type of
Table~\ref{tw:types-table}, and the completions of $K$ at the primes above
$\mathfrak p$ are the local extensions listed there. The type of $\mathfrak p$
in $K/B$ is $(8,4)$ if $\mathfrak p$ lies above $2$, $(2,2)$ if it lies above
$3$ or $5$, and $(1,4)$ if it is a capped prime. In particular $K/F$ has
uniform local types $(e_r,f_r)_{r\in\{2,3,5,29\}}$ in the sense of
Definition~\ref{geo:uniform-types}.
\end{theorem}

\begin{table}[htbp]
\centering
\caption{Absolute types of the places of $K$ and $F$ above $2$, $3$, $5$,
$29$ and above the prime $\mathfrak t_3$ of $B$ above $41$, the corresponding
local extensions of the completions of $B$, and the contributions of these
primes to $\rd(K)$. The last row is the prime $241$, which ramifies in $B$;
$K/B$ is unramified above it.}
\label{tw:types-table}
\begin{tabular}{llll}
\toprule
$r$ & $(e_r,f_r)$ & local extension over $B$ & factor of $\rd(K)$\\
\midrule
$2$ & $(8,4)$ & $L_2/\Q_2$, group $\mathcal D$ & $2^{9/4}$\\
$3$ & $(2,2)$ & $\Q_3(\sqrt{-1},\sqrt3)/\Q_3$, group $C_2\times C_2$ & $3^{1/2}$\\
$5$ & $(2,2)$ & $\Q_5(\sqrt2,\sqrt5)/\Q_5$, group $C_2\times C_2$ & $5^{1/2}$\\
$29$ & $(1,4)$ & unramified of degree $4$ over $\Q_{29}$ & $1$\\
$41$ (above $\mathfrak t_3$) & $(1,4)$ & unramified of degree $4$ over $\Q_{41}$ & $1$\\
\midrule
$241$ & $e=2$ & unramified & $241^{1/2}$\\
\bottomrule
\end{tabular}
\end{table}

\begin{proof}
\emph{Existence.} By Proposition~\ref{tw:infinite}, $G_B$ has open normal
subgroups of arbitrarily large index, and $D_4G_B$ is open because $G_B$ is
finitely generated. Given $m$, choose an open normal subgroup
$\mathcal U'\subset D_4G_B$ of $G_B$ of index at least $m$, and put
$N'=D_4G_B/\mathcal U'$. A nontrivial normal subgroup of a finite $2$-group
meets its center, so $N'$ contains normal subgroups of $G_B/\mathcal U'$ of
every order dividing $|N'|$. Take one of order $[G_B:\mathcal U']/m$, which
divides $|N'|=[G_B:\mathcal U']/2^{49}$. Its preimage $\mathcal U$ is open and
normal in $G_B$, contained in $D_4G_B$, and of index $m$. Let $K$ be its fixed
field in $B_G$. Conversely, every $K$ as in the statement is the fixed field
of an open normal subgroup $\mathcal U\subset D_4G_B$. We fix such a $K$ and put
$H=\Gal(K/B)=G_B/\mathcal U$.

\emph{Complex places.} Since $\sqrt{-1}=\sqrt{\alpha_0}\in E\subset K$, the
field $K$ is totally imaginary, and $[K:F]=2$. Let $j\in\{1,2\}$. The group
$H$ permutes the places of $K$ above $v_j$ transitively. If an embedding
$\sigma\colon K\to\C$ induces such a place and $\iota_\sigma\in H$ is its
complex conjugation, so that $\bar\sigma=\sigma\circ\iota_\sigma$, then for
$h\in H$ we have
$\overline{\sigma\circ h}=(\sigma\circ h)\circ(h^{-1}\iota_\sigma h)$, so the
complex conjugation at the place of $\sigma\circ h$ is
$h^{-1}\iota_\sigma h$. By the definition of
the local map at $v_j$, the image of $\iota_j$ in $H$ is the complex
conjugation at the place of $K$ below the place of $B_S$ fixed in
Definition~\ref{tw:local-generators}. Hence every complex conjugation at a
place of $K$ above $v_j$ is conjugate in $H$ to the image of $\iota_j$, and
has elementary image $\bar\iota_j$, since conjugation preserves elementary
images. In particular, two choices of $\iota_1$ in the statement are
conjugate in $H$, so an element of $H$ maps the field $F$ of one choice onto
that of the other, and the statements of the theorem hold for both choices
or for neither.

Let $\sigma_1\colon K\to\C$
induce the chosen place, so that $\bar\sigma_1=\sigma_1\circ\iota_1$ and
$\sigma_1(F)\subset\R$. The embeddings above $v_1$ are $\sigma_1\circ h$ with
$h\in H$, and two of them agree on $F$ exactly when they differ by right
multiplication by an element of $\langle\iota_1\rangle$. The restriction of
$\sigma_1\circ h$ to $F$ is real exactly when $\sigma_1\circ\iota_1h$ and
$\sigma_1\circ h$ agree on $F$, that is, when
$h^{-1}\iota_1h\in\langle\iota_1\rangle$, that is, when $h$ lies in the
centralizer $C_H(\iota_1)=\{h\in H:h\iota_1=\iota_1h\}$. Similarly, if $\iota'$ is complex conjugation for an
embedding $\sigma_2$ above $v_2$, the restriction of $\sigma_2\circ h$ to $F$
is real exactly when $h^{-1}\iota'h=\iota_1$. This never happens, because
conjugation preserves elementary images and $\bar\iota'=\bar\iota_2\ne\bar\iota_1$
(Lemma~\ref{tw:retained-quotient}(d)). Hence $b=|C_H(\iota_1)|/2\ge1$, so $F$
has a real place and $\theta=(1-b/d)/2<1/2$. Moreover $d=|H|$, so
$b/d=|C_H(\iota_1)|/(2|H|)$ is half the reciprocal of the size $N_\iota$ of
the conjugacy class of $\iota_1$ in $H$, that is, $b/d=1/(2N_\iota)$. This
class maps onto the class in $\overline G_B$ of the element $\iota_1$ of
Definition~\ref{tw:local-generators}, which has $2^{15}$
elements by Lemma~\ref{tw:retained-quotient}(c). Hence $N_\iota\ge2^{15}$,
$b/d\le2^{-16}$ and $\theta\ge(1-2^{-16})/2=\theta_*$.

\emph{Local structure and splitting.} Let $\mathfrak P$ be a prime of $K$
above $\mathfrak p\in S\cup\{\mathfrak t_1,\mathfrak t_2,\mathfrak t_3\}$. Its
decomposition group in $H$ is conjugate to the image of the local map at
$\mathfrak p$. Since $\mathcal U\subset D_3G_B$, Lemma~\ref{tw:retained-quotient}(b)
shows that this image is $\mathcal D$, $C_2\times C_2$ or $C_4$, the last one cyclic on
the Frobenius, and that the prescribed local group injects into $H$.
Therefore the completion $K_{\mathfrak P}$ is $L_2$ over
$B_{\mathfrak p_j}=\Q_2$, the field $\Q_p(\sqrt u,\sqrt p)$ over
$B_{\mathfrak p}=\Q_p$ for $p=3,5$, and the unramified extension of degree
four of $B_{\mathfrak t_k}$. Since $B_{\mathfrak t_1}=B_{\mathfrak t_2}=\Q_{29}$
and $B_{\mathfrak t_3}=\Q_{41}$, the places of $K$ above $\mathfrak p$ have the
absolute types of Table~\ref{tw:types-table}, and the types in $K/B$ stated in
the theorem. Every element of
the decomposition group has elementary image in the span of the images of the
local generators at $\mathfrak p$, whereas every conjugate of $\iota_1$ has
elementary image $\bar\iota_1$. By Lemma~\ref{tw:retained-quotient}(d),
$\iota_1$ lies in no decomposition group of a prime above
$S\cup\{\mathfrak t_1,\mathfrak t_2,\mathfrak t_3\}$, so the decomposition
group of $\mathfrak P$ in $\Gal(K/F)=\langle\iota_1\rangle$ is trivial. Hence
the places of $F$ above $\mathfrak p$ split in $K$, and they have the same
absolute types.

\emph{Ramification and root discriminant.} The extension $K/B$, and hence
$K/F$, is unramified at every prime outside $S$, in particular at
$(\sqrt{241})$. At primes above $S$ the inertia group of $K/F$ is trivial,
because it lies in the decomposition group. So $K/F$ is unramified at every
finite place, its relative discriminant is trivial, $|\Delta_K|=|\Delta_F|^2$,
and $\rd(K)=\rd(F)$. To compute $\rd(K)$ we use
$|\Delta_K|=241^{[K:B]}\mathrm N\mathfrak d_{K/B}$. Since $K/B$ is Galois, for
$\mathfrak p\in S$ all primes $\mathfrak P$ above $\mathfrak p$ have the same
ramification index $e$, residue degree and different exponent
$\operatorname{ord}_{\mathfrak P}(\mathfrak D_{K/B})$, where
$\operatorname{ord}_{\mathfrak P}$ and $\operatorname{ord}_{\mathfrak p}$ are
the normalized valuations at $\mathfrak P$ and $\mathfrak p$, and
$\operatorname{ord}_{\mathfrak p}(\mathfrak d_{K/B})=[K:B]\,\operatorname{ord}_{\mathfrak P}(\mathfrak D_{K/B})/e$.
For $\mathfrak p$ above $3$ or $5$ we have $e=2$ and
$\operatorname{ord}_{\mathfrak P}(\mathfrak D_{K/B})=1$, since over the unramified extension
$\Q_p(\sqrt u)$ the different of $\Q_p(\sqrt u,\sqrt p)$ is generated by
$2\sqrt p$. For $\mathfrak p_1,\mathfrak p_2$, Lemma~\ref{tw:local-two}(d)
gives $\operatorname{ord}_{\mathfrak P}(\mathfrak D_{K/B})/e=9/4$. Since $[K:\Q]=2[K:B]$, the
base contributes $241^{1/2}$, each of the two primes above $3$ contributes
$3^{1/4}$, each of the two above $5$ contributes $5^{1/4}$, and each dyadic
prime contributes $2^{9/8}$. Hence
$\rd(K)=2^{9/4}\,3^{1/2}\,5^{1/2}\,241^{1/2}=\lambda$.
\end{proof}

\begin{remark}\label{tw:same-quotient}
Consider the variant of the construction in which the third cap is placed at
the inert prime $\mathfrak t_0=7\mathcal O_B$ instead of at $\mathfrak t_3$.
Its quotient by the fourth Zassenhaus term, its relation spaces $R_2$ and
$R_3$, and the class of $\iota_1$ are the same as here: the proof of
Lemma~\ref{tw:retained-quotient}(a) uses only the generators of $N_B$ in
(i)--(iii) of Lemma~\ref{tw:GB-presentation} and the commutators
$[[\hat y_j,\hat z_j],\hat z_j]$, since a fourth power lies in $D_4\mathfrak F$.
The two towers also have the same Golod--Shafarevich function, since both
impose three conditions of type $C_4$. They differ at the primes above
$7$ and $41$.
\end{remark}

\begin{remark}\label{tw:lean-remark}
The Lean formalization \cite{NaslundLeanW} formalizes the present tower,
with its caps at $\mathfrak t_1$, $\mathfrak t_2$ and $\mathfrak t_3$.
Remark~\ref{an:lean-hypothesis} states what it proves.
\end{remark}
